\chapter{Study Area, Geological Setting, and Data} \label{chap: Study Area, Geological Setting, and Data}

\section{The Tokachi River Basin: Hydrology, Morphology, and Flood History}
\label{sec: The Tokachi River Basin: Hydrology, Morphology, and Flood History}

The preceding chapter established on theoretical and methodological grounds why the time-dependent progression of backward erosion piping (BEP) is a credible yet historically unquantified threat in high-gradient Japanese river systems. This section grounds that argument in the physical reality of the study basin. It characterises the Tokachi River as a representative flashy system, describes the geomorphology of the Obihiro study reach and the urban floodplain it protects, and reviews the historical flood record that culminates in the August 2016 consecutive typhoon sequence. Together these establish the hydrological loading character, the spatial setting, and the empirical survival evidence that the three-phase probabilistic framework of Chapter~\ref{chap: Methodology} is designed to exploit.

% Suggested figure: basin location and overview map with catchment outline,
% main stem, principal tributaries, flood-control dams, Obihiro,
% and the two study reaches.

\subsection{Hydrological Character: A Flashy River System}
\label{subsec: Hydrological Character: A Flashy River System}

The Tokachi River drains approximately 9{,}010~km\textsuperscript{2} of southeastern Hokkaido with a mainstem length of around 156~km, discharging to the North Pacific \parencite{kimura_2018}. The catchment is bounded by high mountains to the west, north, and east, with elevations exceeding 2{,}000~m in the headwaters, and its planform is unusually compact: the basin shape factor (catchment area divided by the square of mainstem length) is approximately 0.37, exceptionally large among major Japanese rivers \parencite{fukuoka_2019}. This near-circular geometry causes the flood peaks of the principal tributaries to arrive at their confluences almost simultaneously, superimposing into a single sharp mainstem wave rather than a temporally dispersed sequence of sub-peaks \parencite{fukuoka_2019, kimura_2018}. Combined with steep channel gradients that produce a rapid concentration of runoff, the result is a hydrological regime characterised by abrupt water-level rise and comparatively short flood durations, the defining attributes of a flashy system \parencite{kyuka_2020}. It is precisely this brevity of loading that has historically been assumed to protect the basin against duration-driven subsurface failure mechanisms such as BEP, and quantifying the conditions under which that assumption may fail is the central purpose of this thesis.

The principal tributaries relevant to the study reach are the Otofuke River entering from the north and the Satsunai River entering from the south-west, both joining the mainstem in the KP~53 to KP~56 reach; the largest tributary, the Toshibetsu River, enters further downstream near KP~29 \parencite{hoshino2023spatiotemporal, fukuoka_2019}. Runoff from the headwaters reaches the downstream plain with a lag of only several hours \parencite{kimura_2018}, and river planning in the basin adopts a 72-hour rainfall duration because peak discharge generally occurs within that window of the onset of rainfall \parencite{hoshino2023spatiotemporal}. Even for a fixed 72-hour basin-total, the spatial concentration of rainfall within the catchment strongly modulates the resulting hydrograph: \textcite{hoshino2023spatiotemporal} demonstrate that peak discharge at a given gauge can differ by more than a factor of two across ensemble events with equivalent total rainfall volume, depending on which sub-catchment bears the heaviest precipitation. This sensitivity is central to the present framework, which supplies the full continuous multi-peak hydrograph to the time-stepping solver rather than scalar peak values (Section~\ref{sec: The d4PDF Climate Ensemble: Structure, Downscaling, and Hydrograph Extraction}). Flood peaks are further attenuated by the Tokachi Dam on the upper mainstem and the Satsunai River Dam on the Satsunai tributary; adaptive pre-release operation at the Tokachi Dam demonstrably reduced downstream water-stage peaks during the 2016 event \parencite{kimura_2018}.

Within the deterministic Japanese flood-safety framework reviewed in Chapter~\ref{chap: Literature Review}, the basin is managed to a planning-scale return period in the range of 100 to 200 years \parencite{final_report_2022, wp1_report_2021}. The planned high-water discharge is approximately 13{,}700~m\textsuperscript{3}/s at the Moiwa reference gauge (Tokachi KP~21) near the outlet and around 6{,}100~m\textsuperscript{3}/s at the Obihiro reference point (Tokachi KP~57), corresponding to a design annual-maximum 72-hour rainfall of approximately 245.7~mm at that reference point \parencite{fukuoka_2019, wp1_report_2021}. These deterministic design quantities provide the reference loading against which the historical events discussed below, and the climate-perturbed ensemble of Chapter~\ref{chap: Methodology}, are assessed.

\subsection{Geomorphology of the Study Reach and the Urban Floodplain}
\label{subsec: Geomorphology of the Study Reach and the Urban Floodplain}

The geomorphology of the Tokachi basin is the legacy of late-Pleistocene glacial and periglacial processes. Mountain glaciers occupied the Hidaka Range through the last glacial stage, and large volumes of poorly sorted periglacial and fluvial sediment were transported downstream, building extensive alluvial fans and fill-top fluvial terraces across the eastern footslopes and the Tokachi plain \parencite{furuichi_2018}. The relict-channel and overbank deposits laid down by this history are the direct geological origin of the bipartite foundation stratigraphy that governs the piping vulnerability of the study reach: a thin clayey-silt overbank cap (the $A_c$ blanket) overlying coarse fluvial channel-fill gravels (the $A_g$ aquifer), documented in detail in Section~\ref{sec: The Bipartite Stratigraphy: The Ac/Ag Foundation System}. The current-state geomorphological classification of the study reach as overlying former river course and floodplain deposits \parencite{fukuda_2025_internal} is therefore not incidental; it is the surface expression of exactly the depositional environment that produces uplift- and heave-prone layered foundations.

The city of Obihiro occupies the flat alluvial plain in the confluence zone of the Tokachi, Satsunai, and Otofuke Rivers \parencite{kimura_2018}, and the study reach protects this high-value urban floodplain: the Tokachi right bank from KP~53.8 to KP~62.8 and the Satsunai left bank from KP~3.2 to KP~16.6, the spatial boundaries inherited from \textcite{uemura_phd_2025} and detailed in Section~\ref{sec: Spatial Scope: Study Reaches and Segment Discretization}. The Tokachi tributaries were originally multi-thread braided gravel-bed rivers; flow regulation, channelisation, and embankment construction progressively transformed many reaches into single-thread channels with vegetated floodplains \parencite{kyuka_2020}. The embankments protecting the study reach were first raised in 1937 and subsequently enlarged through successive reinforcement campaigns, so that the present cross-sections embody a patchwork of gravelly, sandy, and cohesive fills overlying the relict floodplain substrate \parencite{oyo_1999, fukuda_2025_internal}. A defining geometric attribute of the reach is the pronounced longitudinal variation in foreshore width, spanning from approximately 44~m at KP~62.0 to approximately 600~m at KP~60.0, a more than ten-fold range within nine kilometres \parencite{oyo_1999}. Foreshore width sets the foreland seepage-path length and so enters the analytical head response factor $r_e$ through a finite-width correction. The fragility computed in this work nevertheless proves almost insensitive to it, and the dominant section-to-section control is the thickness of the confining blanket rather than the width of the foreshore (Section~\ref{subsec: Cross-Section Inventory and the Foreshore-Width Control on Risk}).

A geomorphic and geotechnical distinction within the basin must be drawn clearly to avoid conflating two physically different settings. The lower Tokachi reach downstream of the study area is underlain by thick, highly compressible peat that, after long-term consolidation beneath successive embankment additions, attains a near-impermeable state (with hydraulic conductivities on the order of $10^{-11}$ m/s); the governing seepage path there lies within the gently sloping levee (kyuryo-tei) body rather than the foundation \parencite{fukuoka_2019}. This contrasts sharply with the permeable bipartite gravel foundation of the upper Tokachi study reach examined here, a difference that is central to the correct interpretation of the 2016 survival evidence in Section~\ref{sec: The 2016 Consecutive Typhoon Events and Observed Levee Behaviour}. The high-gradient, coarse-bedded character of the basin's rivers also renders them morphodynamically active: during large floods, alternate sandbars propagate and mid-channel bars develop, deflecting flow toward the embankments and driving lateral bank erosion that has produced levee breaches independently of overtopping \parencite{kyuka_2020}. This fluvial-erosion pathway enters the integrated multi-mechanism framework through the pre-calculated scour fragility of \textcite{uemura_phd_2025}, allowing the present work to isolate the contribution of backward erosion piping.

\subsection{Historical Flood Record}
\label{subsec: Historical Flood Record}

As a sub-boreal region, Hokkaido has historically been less frequently affected by intense typhoons than the main island of Japan, and flood-defence design evolved under that assumption \parencite{furuichi_2018}. Notable floods nonetheless occurred: a 28-hour rainfall event in August 1981 delivered 331~mm at Kamisatsunai, then among the largest events in the local record \parencite{furuichi_2018}, and a further large flood in September 2011 drove documented morphological change in the Otofuke tributary and is incorporated as a supplementary survival constraint in Chapter~\ref{chap: Methodology} \parencite{kyuka_2020}. The basin's levees have also been repeatedly tested by seismic loading, with the 1968 and 2003 Tokachi-oki and the 1993 Kushiro-oki earthquakes each causing embankment damage and motivating the progressive cross-section enlargement programme of the lower Tokachi \parencite{fukuoka_2019, oyo_1999}.

The defining event of the modern record is the consecutive typhoon sequence of August 2016. Four systems, Chanthu, Kompasu, Mindulle, and Lionrock, affected Hokkaido within roughly two weeks, an unprecedented clustering in the instrumental record \parencite{kimura_2018}. The three preceding systems saturated basin soils before the arrival of the fourth; a counterfactual simulation by \textcite{kimura_2018} that removes the antecedent rainfall reduces downstream water-stage peaks by approximately 4 to 24~per cent across gauging stations, empirically substantiating the compound-loading mechanism identified in Chapter~\ref{chap: Literature Review} as disproportionately dangerous for seepage-driven failure. Typhoon Lionrock followed an unusual eastward track that drove intense orographic rainfall against the eastern slopes of the mountains, delivering more than 500~mm over three days at the heaviest locations \parencite{kimura_2018, furuichi_2018}; at the Karikachi gauge, the maximum 24-hour (352~mm) and 72-hour (507~mm) totals corresponded to recurrence intervals of roughly 110 and 109~years respectively \parencite{furuichi_2018}. The event caused agricultural damage exceeding 40{,}000~ha and economic losses of around 260~million~USD across Hokkaido \parencite{kimura_2018}, and at the Moiwa reference gauge the reconstructed mainstem discharge reached approximately 10{,}870~m\textsuperscript{3}/s against the design discharge of 13{,}700~m\textsuperscript{3}/s, with water levels exceeding the Planned High Water Level by roughly 1~m along most of the directly managed reach \parencite{fukuoka_2019}.

The pattern of damage was spatially differentiated in a manner that directly motivates the multi-mechanism scope of this thesis. In the steep granitic headwater catchments of the western basin, the dominant processes were sediment-related: \textcite{furuichi_2018} document disastrous debris flows and bank erosion in weakly cohesive periglacial deposits, with total sediment discharge on the order of $5\times10^{6}$~m\textsuperscript{3} to the downstream reaches. In the Otofuke River, four flood peaks within 15 days drove rapid lateral channel migration and meander shift, producing levee breaches at seven locations without overtopping, as the water level remained below the crest throughout \parencite{kyuka_2020}. Along the Satsunai and its Tottabetsu tributary, a levee breach caused approximately 50~ha of inundation \parencite{fukuoka_2019}. These morphodynamic and fluvial-scour failures are represented in the integrated framework through the pre-calculated surface fragility of \textcite{uemura_phd_2025}.

Against this backdrop of severe damage elsewhere, the behaviour of the study levees constitutes the central empirical paradox that the Bayesian reliability updating of Phase~2 is designed to resolve. Severe sand boiling and piping-induced deterioration were observed along the Tokoro River levees during the same 2016 event \parencite{kawajiri_2025, tokorogawa_2017}, confirming that the regional layered stratigraphy is physically vulnerable to seepage. Yet official post-disaster investigation confirmed no sand boils or piping along the Tokachi and Satsunai study segments \parencite{tokachi_levee_committee_2017}, despite the 1998 evaluation having formally rated several cross-sections as seepage-deficient \parencite{oyo_1999}. \textcite{fukuoka_2019} provide a mechanistic basis for this resilience in the wide gently sloping levees of the lower Tokachi: the levee vulnerability index $t^{*}$ remained at approximately $10^{-3}$ throughout the 2016 loading, an order of magnitude below the slope-failure threshold, yet approaches the danger zone under full design-scale discharge. The 2016 record therefore furnishes both the worst-case compound hydrograph and the documented non-failure observation that calibrates the prior fragility curves; the detailed contrast between the piping-affected Tokoro River and the surviving study segments, and its significance for the calibration, is examined in Section~\ref{sec: The 2016 Consecutive Typhoon Events and Observed Levee Behaviour}.

\section{The 2016 Consecutive Typhoon Events and Observed Levee Behaviour}
\label{sec: The 2016 Consecutive Typhoon Events and Observed Levee Behaviour}

The compound typhoon disaster of August 2016 occupies a dual position in this study. It is simultaneously the most severe hydraulic loading event in the modern observational record for the region and, for the specific levee segments examined here, a documented survival. This duality is what makes the event the central empirical reference of the present work. On one hand, its fully resolved hydrograph supplies the loading constraint against which the prior backward erosion piping (BEP) fragility curves are calibrated in the Bayesian updating of Phase~2 (Chapter~\ref{chap: Methodology}). On the other, when the survival of the Tokachi and Satsunai study reaches is read against the contemporaneous seepage failure of the neighbouring Tokoro River levees, it resolves the apparent contradiction between the 1998 deterministic deficiency rating of Table~\ref{tab:oyo_1998} and the observed field resilience. Whereas Section~\ref{sec: The Tokachi River Basin: Hydrology, Morphology, and Flood History} characterised the meteorology and basin-scale hydrology of the event, the present section concentrates on the loading transmitted to the study cross-sections and on the levee behaviour that was actually observed, since it is this behaviour that the calibration procedure exploits as data.

\subsection{The August 2016 Compound Event}
\label{subsec: The August 2016 Compound Event}

The defining feature of the August 2016 disaster, for the purposes of a duration-governed failure mechanism, is its compound, multi-peak character rather than the magnitude of any single peak. As documented in Section~\ref{subsec:flood_history}, four typhoons reached Hokkaido within a 15-day window: Chanthu, Kompasu and Mindulle in succession between 17 and 23 August, followed by the anomalously tracked Lionrock on 30 to 31 August \parencite{kimura_2018}. The three predecessors delivered approximately 60 per cent of the total event precipitation, exhausting the shallow surface storage of the basin before the most intense rainfall arrived \parencite{kimura_2018}. Distributed hydrological simulation confirms the consequence of this sequencing: flood peaks during the approach of the fourth typhoon were between 4 and 24 per cent lower when the antecedent rainfall of the preceding storms was removed from the model, with the amplification strongest in the early stage of the final event and decaying thereafter as fresh rainfall came to dominate the runoff response \parencite{kimura_2018}. On the Tokachi main stem the resulting peak discharge reached approximately 79 per cent of the design high-water discharge at the Moiwa reference section, with the water surface standing roughly 1~m above the planned high-water level through most of the directly managed reach \parencite{fukuoka_2019}. At the study segments themselves, the loading is referenced to the upstream Obihiro gauge rather than to the downstream Moiwa section at KP~20.8 \parencite{oyo_1999}, and lies within the directly managed reach that rose above the planned high-water level in 2016 \parencite{fukuoka_2019}. For the Tokachi right-bank segments the documented high-water-level duration is approximately 34~hours over KP~56 to KP~57 and approximately 24~hours from KP~58 onward, with a mean hydraulic gradient across the levee base in the range 0.1 to 0.2 \parencite{fukuda_2025_internal}. The study segments therefore sustained an elevated river stage for a substantial period rather than only a brief transient peak, so that their complete absence of sand boiling or piping (Section~\ref{subsec: Levee Performance: Tokoro River versus the Tokachi and Satsunai Study Reaches}) cannot be ascribed to trivially weak external loading and is examined instead in terms of foreshore-controlled head transmission and the progression race in Chapter~\ref{chap: Methodology}.

This compound structure is precisely the loading signature that a single-peak design hydrograph, of the kind underlying the 1998 deterministic evaluation (Section~\ref{subsec: Geotechnical Dataset Design Hydrograph}), cannot reproduce. Because pipe progression accumulates irreversibly across successive sub-peaks under the memory formulation of Section~\ref{sec: Compound Event Modelling: Cumulative Memory and Irreversible Pipe Growth}, the partially eroded pipe length carried forward from one peak to the next is a quantity that only a continuous multi-peak record can constrain. For this reason the 2016 event is retained throughout this work as a full transient time series, denoted $h_{2016}(t)$, rather than as a scalar peak; its role as the calibration boundary is developed in Section~\ref{sec: The Historical Constraint: The 2016 Transient Hydrograph as Empirical Boundary}. The well-documented September 2011 flood, of lower peak magnitude, is retained as a supplementary and independent survival anchor.

\subsection{Levee Performance: Tokoro River versus the Tokachi and Satsunai Study Reaches}
\label{subsec: Levee Performance: Tokoro River versus the Tokachi and Satsunai Study Reaches}

The evidential value of the 2016 survival rests on a contrast between two Hokkaido river systems that experienced the same compound event with comparable layered foundations but markedly different seepage outcomes. The Tokoro River is an entirely separate river system in eastern Hokkaido, draining a catchment of approximately 1{,}930~km\textsuperscript{2} over a trunk stream of some 120~km to the Sea of Okhotsk, passing the city of Kitami, and shares the steep, flashy hydrological character of eastern Hokkaido rivers \parencite{tokorogawa_2017}. In its affected downstream reach the channel narrows to a width of only 100 to 200~m as it threads a constricted mountain pass, a geometry that promoted prolonged elevation and retention of the flood stage at the levees \parencite{morita_2018}. The loading there was both severe and sustained: at the Futochanae observatory the river stage exceeded the design high-water level on multiple occasions across the compound sequence and, during the most intense peak, remained continuously above that level for a period on the order of tens of hours \parencite{morita_2018}. This is a directly documented instance, in a neighbouring flashy Hokkaido system, of the extended-duration multi-peak loading that the present framework treats as the governing condition for retrogressive pipe development.

During the August 2016 event the Tokoro levee system sustained extensive seepage-related distress along the KP16 to KP27 reach: numerous sand boils and concentrated leakage were recorded, two tributary levees breached, including a 100~m breach of the Shibayamazawa levee, and overflow with associated landside slope failure occurred at several sections, producing inundation of roughly 500~ha \parencite{tokorogawa_2017}. The boils were concentrated within an approximately 4~km span of the affected reach, and no leakage or sand boiling had been recorded at these locations during previous floods \parencite{morita_2018}.

Critically, the post-disaster survey committee distinguished the seepage-driven distress from the overflow-driven breaches through site-specific seepage analyses informed by the observed 2016 water levels. At the worst-affected sand-boiling location, near KP26.8 on the Tokoro, the computed vertical and horizontal exit gradients of 0.70 and 0.87 both exceeded the Japanese critical threshold of $i_c = 0.5$, whereas at the overflow-affected sections farther upstream, near KP18.9 and KP22.6, the computed exit gradients all remained subcritical, in the range 0.24 to 0.45 \parencite{tokorogawa_2017}. The grain-size composition of the ejected material matched the shallow sandy foundation layer rather than the embankment fill \parencite{kawajiri_2025}, and detailed geotechnical investigation of the boiling reach confirmed that this sandy source layer was confined beneath a low-permeability silt-clay surface cap, logged in the regional cohesive, sand and sand-gravel ($A_c$/$A_s$/$A_g$) nomenclature that is consistent with the bipartite stratigraphy of the Tokachi study reach \parencite{morita_2018}. Taken together, these observations establish that the layered alluvial foundations characteristic of the region are physically capable of seepage-driven initiation under 2016-class loading, and that the sand boiling on the Tokoro was a genuine seepage phenomenon rather than a secondary effect of overtopping.

The Tokoro field investigations additionally clarify the role of foundation structure in governing the seepage response, in a manner directly relevant to the leakage-length argument of Section~\ref{subsec: Stratigraphic Control on the Leakage Length}. Two distinct boil-occurrence patterns were resolved \parencite{kawajiri_2025, morita_2018}: where the sandy source layer was continuous from the riverside to the landside, the largest boils formed directly at the landside toe, the point of easiest pressure release; whereas where a thicker landside silt layer interrupted this path and acted as a dead-end low-permeability barrier, the boils instead emerged at substantial landward distances from the toe, the seepage pressure having been deflected upward by the barrier. In the latter case, multi-channel surface-wave surveys and dynamic cone penetration testing revealed a network of loosened zones rather than a single continuous channel \parencite{kawajiri_2025}. That distal boils could form well landward of the toe is itself consistent with the large leakage length predicted for an $A_c$/$A_g$ stratigraphy, providing field corroboration that elevated aquifer pressures act substantially landward of the levee toe in these layered foundations. It is nonetheless important to note that these post-disaster surveys characterised conditions following boil formation and documented seepage-driven initiation and partial foundation disturbance rather than a completed retrogressive breach traversing the full seepage length \parencite{kawajiri_2025}, the very distinction between initiation and through-breach that the time-dependent framework of this thesis is designed to resolve.

Against this benchmark, the behaviour of the study reaches is the more remarkable. Official post-disaster investigation of the Tokachi right bank and the Satsunai left bank confirmed the complete absence of sand boils and piping during the 2016 event \parencite{tokachi_levee_committee_2017}, despite the 1998 evaluation having rated every one of the five surveyed cross-sections as deficient in uplift and several as deficient in exit gradient, with values reaching $i_v = 0.97$ at KP62.0 on the Tokachi right bank (Table~\ref{tab:oyo_1998}). One interpretation, given physical grounding in Section~\ref{subsec: Historical Field Evidence: The Paradox of Survival and the Tokoro--Tokachi Contrast}, is that this survival reflects the time-dependent nature of the mechanism: the flashy flood wave receded before retrogressive pipe development could traverse the full seepage length, potentially reinforced by sediment-clearing arrest during recession. The Tokoro contrast lends this interpretation substantial weight because it removes one of the competing explanations, namely that the regional stratigraphy is simply not erodible. Seepage-driven initiation is demonstrably possible in Hokkaido's layered foundations and was extensively realised on the Tokoro under sustained 2016-class loading, so the survival of the Tokachi and Satsunai segments cannot be attributed to an absent hazard at the regional scale. The Tokoro contrast does not, however, resolve the matter for the study sections specifically. Several explanations for their survival remain viable and are not eliminated by the occurrence of erosion elsewhere: the section-specific foreland and hinterland geometry entering the response factor $r_e$ (Section~\ref{subsec: Cross-Section Inventory and the Foreshore-Width Control on Risk}) may have attenuated the hydraulic head reaching the foundation, although the foreshore width specifically is shown there to be immaterial; the 1999--2003 side-berm and toe-drain remediation (Section~\ref{subsec: Remediation History and Segment-State Stratification}) may have already neutralised the hazard at the affected nodes; the coarse AgA$_g$ grading may confer static resistance under 2016-class heads; and the critical initiation head may simply not have been reached at the surviving sections. The data assembled in this chapter, namely the section-specific foreshore geometry, the documented remediation states, and the grain-size record, are precisely what permit these alternatives to be separated from the time-dependent explanation within the probabilistic framework of Chapter~\ref{chap: Methodology}. The survival is therefore documented here as an observation to be explained, and the question of which mechanism explains it is intentionally left open for that framework to interrogate rather than answered in advance.

\subsection{Significance for Bayesian Reliability Updating}
\label{subsec: Significance for Bayesian Reliability Updating}

The juxtaposition of a formal deterministic deficiency rating \parencite{oyo_1999} with a documented survival of the most extreme compound flood in the regional record \parencite{tokachi_levee_committee_2017} is exactly the type of informative observation that Bayesian reliability updating is designed to exploit \parencite{schweckendiek_2014, hkv_2023}. A survived water level constitutes evidence $\epsilon$ that the resistance of the levee, expressed as a critical water level $H_c$, exceeded the observed level $h_{obs}$. Conditioning the failure probability on this evidence yields
\begin{equation}
P(F \mid \epsilon) = \frac{P\!\left(H_c < H,\; H_c > h_{obs}\right)}{P\!\left(H_c > h_{obs}\right)},
\end{equation}
where $H$ is the stochastic load. As shown by \textcite{hkv_2023}, this updating admits a closed-form expression for the posterior fragility curve, namely the cumulative distribution $F_{H_c^{*}}$ of the updated critical level,
\begin{equation}
F_{H_c^{*}}(h_c) =
\begin{cases}
\dfrac{F_{H_c}(h_c) - F_{H_c}(h_{obs})}{1 - F_{H_c}(h_{obs})}, & h_c \geq h_{obs}\\[8pt]
0, & h_c < h_{obs}
\end{cases}
\end{equation}
which depends on the survived level $h_{obs}$ but not on its exceedance frequency. Because the updated fragility lies at or below the prior at every water level, the survival-conditioned failure probability never exceeds the prior, and the peak-based closed-form approximation employed in standard WBI2017 practice is provably conservative, satisfying $P(F \mid \epsilon) \leq P_A(F \mid \epsilon) \leq P(F)$ \parencite{hkv_2023, schweckendiek_2016}.

For the time-dependent mechanism treated here, however, this closed-form formulation carries an essential limitation: it conditions only on the scalar peak observed head $h_{\mathrm{obs}}$ and therefore cannot represent the load duration that governs pipe progression. A peak-based survival constraint would credit the levee on the basis of peak stage alone and could falsely exonerate a future long-duration event whose peak lies below that of 2016 but whose extended loading drives sufficient cumulative progression to breach. The present work consequently conditions on the full transient hydrograph $h_{2016}(t)$ through Monte Carlo Accept-Reject filtering, the exact benchmark updating method of \textcite{schweckendiek_2014}, as formalised in Section~\ref{sec: Accept-Reject Filtering: Procedure and Posterior Parameter Distributions}. This is the operational reason the 2016 event is documented in this chapter as a complete time series rather than a peak value: it is the data object upon which the calibration of Chapter~\ref{chap: Methodology} directly operates, trimming from the prior precisely those high-conductivity, thin-blanket parameter combinations that would have predicted a breach the field did not record. The diagnostic strength of this trimming is, however, conditional. The filter is informative about the progression mechanism only insofar as the transient limit state was the binding one at the study sections in 2016; where survival is instead secured by the section's foreshore geometry, entering the model through the section-specific response factor $r_e$, or by a remediation state that suppresses the initiation prerequisites, a realisation survives the filter for reasons unrelated to progression time, and the corresponding reduction in epistemic uncertainty about the time-dependent mechanism is correspondingly small. The framework is constructed to expose exactly this distinction: because the static Sellmeijer limit state is evaluated in parallel with the transient one on the identical parameter sample (Chapter~\ref{chap: Methodology}), the share of the 2016 survival already accounted for by the static, geometry-attenuated, remediation-adjusted limit state can be read directly, and the \emph{marginal} informativeness attributable to the time constraint isolated from it. The calibration should therefore be interpreted as conditional evidence about progression, weighted by the extent to which the simpler explanations have already accounted for the survival, rather than as unconditional confirmation that time-dependence governed the 2016 outcome.

\section{Spatial Scope: Study Reaches and Segment Discretization}
\label{sec: Spatial Scope: Study Reaches and Segment Discretization}

\subsection{Study Reaches and the Inherited Discretization Framework}
\label{subsec: Study Reaches and the Inherited Discretization Framework}

The spatial scope of this study is confined to the river levees protecting the urban centre of Obihiro, comprising the right bank of the Tokachi River over the reach KP~53.8 to KP~62.8 and the left bank of the Satsunai River over the reach KP~3.2 to KP~16.6, where KP denotes the kilometre post measured upstream from the river mouth. These boundaries are not chosen independently; they are adopted verbatim from the antecedent probabilistic framework of \textcite{uemura_phd_2025} and \textcite{uemura_wp2_2024} so that the time-dependent backward erosion piping (BEP) fragility curves developed in this thesis can be integrated, at the identical spatial nodes, with the pre-computed overflow and fluvial scour fragility curves from that work. Maintaining this one-to-one spatial correspondence is both a scientific and a computational prerequisite for the series-system integration described in Chapter~\ref{chap: Methodology}.

Following the discretization methodology of \textcite{uemura_wp2_2024}, the levees within these boundaries are evaluated at 200~metre intervals, which constitute the foundational evaluation segments of the model. These segments are subsequently aggregated into larger sections (Sections~1 to~5 for the Tokachi River and Sections~1 to~4 for the Satsunai River) on the basis of shared inundation characteristics, each section being represented by the segment used to compute its flood consequences \parencite{uemura_wp2_2024}. Because BEP is a weakest-link mechanism, the conditional failure probability of a finite-length segment exceeds that of any single representative cross-section, and the per-cross-section probabilities derived from the Monte Carlo engine are subsequently scaled to the 200~metre segment level using the spatial autocorrelation of the governing parameters, as developed in Section~\ref{sec: The Length Effect and Spatial Autocorrelation of Governing Parameters} following \textcite{kanning_2012} and \textcite{hoffmans_2014}.

The geotechnical foundation for the Tokachi right bank is drawn from a detailed levee strengthening investigation executed by OYO Corporation for the Obihiro Development and Construction Department of the Hokkaido Regional Development Bureau \parencite{oyo_1999}. That investigation furnishes continuous borehole logs, laboratory soil tests, field permeability tests, and a formal seepage and slope safety evaluation along the survey strip KP~56.0 to KP~66.2, organised into five evaluation intervals. Five fully worked-up cross-sections are supplied, at KP~57.4, KP~58.8, KP~60.0, KP~62.0, and KP~63.4 (the last appearing as KP~62.80 in the longitudinal sheet and as KP~63.40 in the borehole panels). These five cross-sections are mapped one-to-one onto the 200~metre evaluation nodes of the study reach and provide the empirical anchor for the prior parameter distributions of Section~\ref{sec: Prior Parameter Distributions: Derivation and Statistical Fitting}.

Two reaches lie outside the borehole coverage of this dataset. The lower Tokachi reach (KP~53.8 to KP~56.0) and the entire Satsunai left bank (KP~3.2 to KP~16.6) are not surveyed in the secured OYO dataset. For these segments, geotechnical parameters are interpolated from the nearest documented cross-sections and assigned deliberately inflated prior uncertainty, consistent with the scope assumption introduced in Chapter~\ref{chap: Introduction}, with supplementary data anticipated from the Obihiro Development and Construction Department through the Japanese research partners \parencite{fukuda_2025_internal}. The numerical priors tabulated in Section~\ref{sec: Prior Parameter Distributions: Derivation and Statistical Fitting} are therefore strictly representative of the five OYO cross-sections, and the interpolation and inflation procedure for the borehole-free segments is treated as a distinct source of epistemic uncertainty.

\subsection{Cross-Section Inventory and the Section-Scale Controls on Risk}
\label{subsec: Cross-Section Inventory and the Foreshore-Width Control on Risk}

The principal geometric attributes of the five OYO cross-sections are summarised in Table~\ref{tab:oyo_inventory}. Each section is resolved by three boreholes, typically one on the landside of the toe, one near the crest, and one on the riverside, which together constrain the contact between the confining blanket and the underlying aquifer and the position of the landside seepage exit.

\begin{table}[htbp]
\centering
\footnotesize
\caption{Cross-section inventory for the five worked-up OYO sections on the Tokachi right bank \parencite{oyo_1999}. Ground elevation gives the range across the three boreholes of each section.}
\label{tab:oyo_inventory}
\begin{tabular}{lccccc}
\toprule
Eval. KP & Boreholes & Ground elev. (m) & HWL (m) & Foreshore width (m) & Foreshore crest level (m) \\
\midrule
57.40 & 3 & $+40.5$ to $+42.5$ & 39.51 & 200 & $+37.98$ \\
58.80 & 3 & $+42.0$ to $+44.1$ & 41.33 & 325 & $+39.10$ \\
60.00 & 3 & $+43.6$ to $+45.2$ & 43.06 & 600 & $+41.60$ \\
62.00 & 3 & $+46.6$ to $+48.8$ & 46.68 & 44  & $+45.80$ \\
63.40 & 3 & $+48.5$ to $+51.1$ & $\approx 49.0$\textsuperscript{a} & river-tight & at grade \\
\bottomrule
\end{tabular}

\vspace{0.5em}
{\footnotesize \textsuperscript{a}~The KP~63.4 high-water level is referenced to the Shikaribetsu River confluence gauge, not the Obihiro gauge, and is derived from the section's own loading basis (initial water level $45.97$~m, with a flood peak of approximately $49.0$~m). This value differs from the $46.68$~m of KP~62.0, with which it should not be conflated.}
\end{table}

The single most consequential geometric attribute for the spatial distribution of piping risk is the foreshore width, which spans from 44~metres at KP~62.0 to 600~metres at KP~60.0, a fourteen-fold variation within a reach of only nine kilometres. It is important to be precise about what this variation does and does not control, because the recorded quantity is easily misread. It is the Japanese \emph{k\=osuishiki-haba}, the width of the high-water bed: the vegetated terrace standing one step above the low-water channel, dry in normal flow and inundated only during events. It is not the distance from the levee to the waterline, which on this braided gravel-bed reach is several hundred metres of unblanketed gravel bar contributing no entry resistance. So defined, the foreshore width is precisely the riverside blanket length $L_1$ of blanket theory \parencite{USACE2000}, and it enters the model through the finite-width correction $\lambda_\mathrm{out,eff} = \lambda_\mathrm{out}\tanh(B_f/\lambda_\mathrm{out})$ on the entry leakage length, and thence the response factor $r_e$.

Two consequences follow, both running counter to the intuition that a narrow foreshore is straightforwardly the dangerous case. First, the correction is saturated at every production section: the realised $\tanh$ credits are $0.969$, $0.995$, $1.000$ and $0.835$ at KP~57.4, KP~58.8, KP~60.0 and KP~62.0, so above roughly $2.5\lambda_\mathrm{out}$ the term becomes numerically indistinguishable from its semi-infinite limit and the 600~metre foreshore earns no more credit than a 100~metre one would. Second, and more consequentially, KP~62.0 does not transmit the most head despite having the narrowest foreshore: its response factor at the prior means is $r_e = 0.330$, the \emph{lowest} of the four sections, against $0.438$, $0.436$ and $0.417$ at KP~57.4, KP~58.8 and KP~60.0, because the same thin and relatively conductive blanket that minimises its uplift resistance also shortens the hinterland leakage length $\lambda_\mathrm{in}$ standing in the numerator. The narrow foreshore raises $r_e$ by only about six per cent relative to an infinitely wide one, which is nowhere near enough to overturn that ordering.

The 1998 exit-gradient contrast between KP~62.0 ($i_v = 0.97$) and KP~57.4 ($i_v = 0.04$) should therefore not be read as a demonstration of foreshore control. Those values are outputs of OYO's own finite-element seepage model, which schematised a continuous cohesive layer across the entire domain \emph{including} the foreshore, so $i_v = 0.97$ is what a fully blanketed model yields and is not evidence of an open foreland \parencite{oyo_1999}. What the shared design hydrograph of Section~\ref{subsec: Geotechnical Dataset Design Hydrograph} does establish is the weaker but still useful point that the section-to-section contrast is geometric and stratigraphic rather than an artefact of differing loading. The section-specific treatment of $r_e$ carried through Chapter~\ref{chap: Methodology} \parencite{pol_thesis_2022, pol_sie_2024} is retained on that basis, while the 200~metre discretization is inherited from \textcite{uemura_phd_2025} rather than derived from foreshore geometry.

\subsection{Longitudinal Characterisation of the Study Reach}
\label{subsec: Longitudinal Characterisation of the Study Reach}

The current-state longitudinal documentation for KP~56 to KP~64 \parencite{fukuda_2025_internal} characterises the reach as a single continuous embankment segment with a mean landside slope of approximately 1:3 to 1:4. The observed and design high-water-level duration is approximately 34~hours over KP~56 to KP~57, stepping to approximately 24~hours from KP~58 onward, and the mean hydraulic gradient across the levee base lies in the range 0.1 to 0.2. The original embankment was constructed in 1937, with an embankment soil that is a patchwork of gravelly, sandy, and cohesive fills. The flood-control geomorphological classification identifies the reach as overlying former river course and floodplain deposits, a depositional setting that is directly consistent with the bipartite stratigraphy documented in Section~\ref{sec: The Bipartite Stratigraphy: The Ac/Ag Foundation System}, in which a thin clayey cap overlies coarse relict-channel gravels. Several hydraulic structures interrupt the reach, including the Tokachi River Bridge, the Kinohara sluice gate (1976), the Heigen Bridge, the Fushiko sluice gate (1980), and the Nakajima Bridge; the influence of structure-adjacent transitions on local seepage is not modelled explicitly, consistent with the technical scope of Chapter~\ref{chap: Introduction}.

\subsection{The 1998 Deterministic Safety Evaluation}
\label{subsec: The 1998 Deterministic Safety Evaluation}

A formal seepage and slope safety evaluation of the study cross-sections was conducted in 1998 under the predecessor framework to the 2002 national levee design guidelines \parencite{oyo_1999}. The results, reproduced in Table~\ref{tab:oyo_1998}, constitute the documented geotechnical motivation for the present study. The evaluation assessed structural adequacy against two governing criteria: whether the local exit hydraulic gradient at the landside toe exceeded the critical threshold $i_c \geq 0.5$, and whether the landside slope stability safety factor fell below the acceptable range of approximately 1.1 to 1.2.

\begin{table}[htbp]
\centering
\footnotesize
\caption{1998 deterministic safety evaluation of the five study cross-sections \parencite{oyo_1999}. Bold entries denote failing criteria. The vertical exit gradient $i_v$ is compared against $i_c \geq 0.5$; the landside slope factor of safety is compared against the acceptable range of approximately 1.1 to 1.2.}
\label{tab:oyo_1998}
\begin{tabular}{lccccc}
\toprule
KP & Riverside slope $F_s$ & Vertical exit $i_v$ & Horizontal $i_h$ & Uplift & Landside slope $F_s$ \\
\midrule
57.40 & 1.229 & 0.040          & 0.050          & \textbf{fail} & \textbf{1.104} \\
58.80 & 1.476 & \textbf{1.300} & \textbf{0.620} & \textbf{fail} & \textbf{1.148} \\
60.00 & 1.448 & \textbf{0.500} & 0.400          & \textbf{fail} & \textbf{1.076} \\
62.00 & 1.375 & \textbf{0.970} & \textbf{0.660} & \textbf{fail} & \textbf{1.172} \\
63.40 & 1.257 & 0.280          & 0.450          & \textbf{fail} & \textbf{1.103} \\
\bottomrule
\end{tabular}
\end{table}

All five sections fail the landside slope stability criterion; KP~58.8, KP~60.0, and KP~62.0 additionally fail the local gradient check; and uplift is flagged as failing across all five sections. Notably, KP~57.4 is the one section to pass the gradient check ($i_v = 0.04$), and the gradient-failing sections are those with narrower foreshores. The correlation is real but must not be read causally: in the model developed here, removing KP~57.4's 200~metre foreshore in its entirety changes its computed failure probability by about $10^{-3}$, so the foreshore is not what secures that pass. Under a common loading the contrast is geometric and stratigraphic rather than an artefact of forcing, but the operative geometry is the confining blanket and the seepage length, not the foreshore width (Section~\ref{subsec: Cross-Section Inventory and the Foreshore-Width Control on Risk}). The collective deficiency rating recorded in these evaluations is the documented basis for the formal deficiency classification of the reach \parencite{fukuda_2025_internal} and is broadly representative of a systemic national vulnerability, given that nationwide seepage inspections subsequently rated approximately 30~per cent of Japan's nationally managed levees as structurally deficient against seepage \parencite{pwri_2013}.

This historical evaluation is decisive for the present framework for two reasons. First, it confirms that the physical prerequisites for piping were not merely present but had been formally quantified as an active engineering threat prior to the consecutive typhoons of August 2016. Second, the apparent paradox between this deterministic warning and the documented survival of these exact segments during the 2016 event, with no sand boils or piping observed \parencite{tokachi_levee_committee_2017}, is precisely the empirical constraint that the Bayesian reliability updating of Chapter~\ref{chap: Methodology} is designed to resolve.

\subsection{Remediation History and Segment-State Stratification}
\label{subsec: Remediation History and Segment-State Stratification}

Following the 1998 deficiency rating, a remediation programme was implemented along affected sections of the reach between 1999 and 2003 \parencite{fukuda_2025_internal}. The works comprised side-berm widening and the installation of landside toe drains along the North Obihiro and East Obihiro levees on the Tokachi right bank. The resulting coexistence of unreinforced sections, berm-widened sections, and drain-equipped sections within a single study reach constitutes a natural spatial experiment in risk stratification, and the model parameterization must reflect the current intervention type at each 200~metre node.

Three intervention states are distinguished, each with a distinct treatment in the limit-state framework of Chapter~\ref{chap: Methodology}. For drain-equipped segments, a functioning landside toe drain actively relieves seepage pressure at the exit point; this is implemented by setting the exit head boundary to zero, under which the uplift and heave limit states remain satisfied for all realisable flood stages, the erosion indicator $I_{er}(t)$ never activates, and the resulting near-zero BEP probability is a physically accurate representation of the drain function rather than a conservative model exclusion \parencite{ENW2017}. For berm-only segments, the post-remediation geometry is applied through the increased seepage length $L$, while the original blanket and aquifer foundation parameters are retained, since the surface earthworks do not alter the underlying stratigraphy. For unreinforced segments, the original geometry is retained.

Two practical limitations attend this stratification. The transmissivity of the post-1999 side-berm fill is recorded as difficult to obtain, with only construction quality-control density data likely to survive, so berm fill is not treated as an independent stochastic conductivity. More fundamentally, the allocation of drain-equipped, berm-only, and unreinforced states to individual 200~metre nodes, and the post-remediation seepage length $L$ for any reinforced segment, must be established from the current-state cross-sections published by the development bureau, which are not contained in the OYO dataset \parencite{fukuda_2025_internal}. The seepage length used for any reinforced segment is therefore taken from the post-remediation geometry rather than from the 1998 OYO sections.


\section{The Bipartite Stratigraphy: The \texorpdfstring{$A_c$/$A_g$}{Ac/Ag} Foundation System}
\label{sec: The Bipartite Stratigraphy: The Ac/Ag Foundation System}

\subsection{Stratigraphic Nomenclature}
\label{subsec: Stratigraphic Nomenclature}

The foundation stratigraphy of the Tokachi right bank, as logged by \textcite{oyo_1999}, conforms closely to the idealised bipartite configuration that predisposes a levee to the composite STPH failure pathway reviewed in Chapter~\ref{chap: Literature Review}. The principal stratigraphic units and their roles in the present model are summarised in Table~\ref{tab:strat_nomenclature}.

\begin{table}[htbp]
\centering
\footnotesize
\caption{Stratigraphic nomenclature of the Tokachi right bank foundation, after the OYO geological legend \parencite{oyo_1999}, with the role assigned to each unit in the present BEP framework.}
\label{tab:strat_nomenclature}
\begin{tabular}{llll}
\toprule
Symbol & Age & Material & Role in BEP model \\
\midrule
$B_c$ & Holocene & silt, gravel-mixed silt (fill)            & embankment \\
$B_s$ & Holocene & silt-mixed fine sand (fill)               & embankment \\
$B_g$ & Holocene & sand-gravel (fill)                        & embankment \\
$A_c$ & Holocene & silt to sandy silt (floodplain deposit)   & confining blanket \\
$A_g$ & Holocene & sand-gravel (floodplain deposit)          & aquifer (primary) \\
$N_s$ & Holocene & sand-gravel (Nakasatsunai terrace)        & aquifer (lower) \\
$K_{1p}$ & Pleistocene & volcanic ash (Kamisatsunai terrace) & effective impervious base \\
\bottomrule
\end{tabular}
\end{table}

The thesis bipartite system is therefore the confining blanket $A_c$ overlying the aquifer $A_g$, supplemented at depth by the coarse $N_s$ terrace gravels, with the assemblage bounded below by the low-permeability volcanic-ash unit $K_{1p}$, which acts as the effective impervious base of the seepage domain. The depositional origin of this configuration is recorded directly in the flood-control geomorphological classification of the reach, namely former river course and floodplain deposits \parencite{fukuda_2025_internal}: the thin clayey $A_c$ cap is the overbank fines deposited atop the coarse fluvial $A_g$ channel fill, which is the classic uplift- and heave-prone arrangement. This regional configuration is mechanistically analogous to the layered stratigraphy that contributed to documented piping failures elsewhere in Japan and to the severe sand boiling observed along the Tokoro River during the 2016 extreme events \parencite{kawajiri_2025, tokorogawa_2017}.

\subsection{Per-Section Stratigraphy and Model-Layer Thicknesses}
\label{subsec: Per-Section Stratigraphy and Model-Layer Thicknesses}

The blanket and aquifer thicknesses extracted from the borehole panels and the corresponding analysis cross-sections \parencite{oyo_1999} are summarised in Table~\ref{tab:strat_thickness}. The aquifer thickness $D_\mathrm{aq}$ is taken as the combined $A_g$ and $N_s$ thickness down to the $K_{1p}$ base, and the blanket thickness $D_{bl}$ as the $A_c$ thickness at the landside toe.

\begin{table}[htbp]
\centering
\footnotesize
\caption{Per-section stratigraphy and model-layer thicknesses on the Tokachi right bank \parencite{oyo_1999}. $D_{bl}$ is the confining $A_c$ blanket thickness; $D_\mathrm{aq}$ is the combined $A_g$/$N_s$ aquifer thickness to the $K_{1p}$ base.}
\label{tab:strat_thickness}
\begin{tabular}{llccl}
\toprule
KP & Embankment & $D_{bl}$ ($A_c$, m) & $D_\mathrm{aq}$ ($A_g$/$N_s$, m) & Notes \\
\midrule
57.40 & $B_c$ (silty)        & 0.80 & 7  & $A_c$ thin homogeneous silt, $\approx$0.8~m ($N=3$); $N>30$ in $A_g$, frequently $>50$ in $N_s$; base $\approx$ EL~23 \\
58.80 & $B_c + B_s$          & 0.85 & 8  & $A_c$ thin sandy silt, $\approx$0.85~m ($N=5$); thick $A_g$ \\
60.00 & $B_s + B_g$          & 0.85 & 9  & $A_c$ thin sandy silt, 0.85 to 1.35~m, thickening riverside (landside toe $\approx$0.85~m, $N=5$ to 6); fine sandy upper $A_g$ over gravel; base $\approx$ EL~30 \\
62.00 & $B_g$                & 0.45 & 10 & $A_c$ thin/transitional, 0.3--0.6~m; high blow-count $A_g$; 44~m foreshore; base $\approx$ EL~32 \\
63.40 & $B_g$                & 1.0 & 11 & $A_c$ effectively absent; modelled as a single gravelly unit \\
\bottomrule
\end{tabular}
\end{table}

The confining blanket is uniformly thin across all four confined sections, spanning only 0.45 to 0.85~m, so blanket-controlled vulnerability is governed by where within this narrow band each section falls rather than by any genuinely thick aquitard. The thinnest blanket occurs at KP~62.0, where the $A_c$ unit is a thin, gravelly, transitional layer of only 0.3 to 0.6~m at its updip pinch-out; the thickest, at KP~58.8 and KP~60.0, reaches only about 0.85~m at the landside toe, with KP~60.0 thickening to 1.35~m toward the riverside but thinning to its landside value at the exit where $D_{bl}$ is evaluated. KP~57.4 sits just below these at about 0.8~m. A short distance upstream at KP~63.4 the $A_c$ is absent altogether and the OYO model collapses the foundation into a single gravelly unit ($D_{bl} \approx 1$~m, nominal). Because both the uplift limit state and the heave gradient scale inversely with blanket thickness, KP~62.0's 0.45~m blanket gives it the worst-case configuration for initiation among the confined sections, consistent with its failing 1998 uplift and gradient ratings and its role as the governing piping section. It is worth being explicit that the blanket thickness alone accounts for this, and that it cuts in both directions: the same thin, conductive $A_c$ that minimises the uplift and heave resistance also shortens $\lambda_\mathrm{in}$ and therefore gives KP~62.0 the lowest response factor of the four sections ($r_e = 0.330$), so it transmits the \emph{least} head to its toe rather than the most. Its narrow 44~m foreshore contributes only the smallest $\tanh$ credit, $0.835$ against $0.969$ to $1.000$ elsewhere, worth about six per cent on $r_e$ and immaterial to the computed fragility (Section~\ref{subsec: Cross-Section Inventory and the Foreshore-Width Control on Risk}). KP~63.4 is in any case anomalous in several other respects, including its use of a different design hydrograph, an aquifer conductivity roughly one and a half orders of magnitude below the other sections, and a distinct borehole naming convention, and it is consequently treated as a structurally distinct section throughout the prior derivation of Section~\ref{sec: Prior Parameter Distributions: Derivation and Statistical Fitting}, with a deliberate decision required as to whether it belongs in the same fragility population as the remaining sections.

\subsection{Stratigraphic Control on the Leakage Length}
\label{subsec: Stratigraphic Control on the Leakage Length}

The hydraulic consequence of the bipartite stratigraphy is a large leakage length. With a thin, low-conductivity $A_c$ blanket confining a thick, highly conductive $A_g$ aquifer, the hinterland leakage length
\begin{equation}
\lambda = \sqrt{\frac{k_\mathrm{aquifer}\cdot D_\mathrm{aq}\cdot D_{bl}}{k_{bl}}}
\end{equation}
is large \parencite{bezuijen_2017}, so that elevated aquifer pressures act substantially landward of the levee toe rather than being confined to its immediate vicinity. This theoretical prediction is consistent with the distal sand boil locations documented along the Tokoro River during the 2016 events \parencite{kawajiri_2025}, and it provides a field-corroborated basis for the spatial extent of the modelled pressure influence. The same leakage length enters the analytical head response factor $r_e$ that translates river stage into aquifer overpressure (Chapter~\ref{chap: Methodology}); combined with the foreshore-width variation of Section~\ref{subsec: Cross-Section Inventory and the Foreshore-Width Control on Risk}, this renders $r_e$, and hence the initiation prerequisites, strongly dependent on the section-specific stratigraphy. The blanket and aquifer conductivities that govern $\lambda$ are accordingly promoted to stochastic variables in Section~\ref{sec: Prior Parameter Distributions: Derivation and Statistical Fitting}, rather than being held at fixed values.

\section{Geotechnical Dataset: Borehole Logs and Laboratory Testing}
\label{sec: Geotechnical Dataset: Borehole Logs and Laboratory Testing}

\subsection{Investigation Campaign}
\label{subsec: Investigation Campaign}

The geotechnical dataset underpinning the prior distributions was acquired during a levee strengthening investigation commissioned for fiscal year 1998 by the Obihiro Development and Construction Department of the Hokkaido Regional Development Bureau and executed by OYO Corporation \parencite{oyo_1999}. The survey strip extends along the Tokachi right bank from KP~56.0 to KP~66.2, organised into five evaluation intervals, and directly covers the majority of the present study reach. The transverse layout at each evaluation section typically comprises one borehole landside of the toe, one near the crest, and one on the riverside, a configuration that resolves both the contact between the confining blanket and the underlying aquifer and the position of the landside seepage exit at which the initiation limit states are evaluated. The campaign combined cored boreholes, recorded with both longitudinal and transverse coverage symbols, with supplementary Swedish-type soundings.

\subsection{Laboratory and Field Test Programme}
\label{subsec: Laboratory and Field Test Programme}

The laboratory programme reported per borehole comprises grain-size analysis, consistency (Atterberg) testing, constant-head permeability testing on density-adjusted reconstituted specimens, and consolidated-undrained triaxial shear testing, again on density-adjusted specimens. The grain-size analysis reports the gravel, sand, silt, and clay percentages, the maximum particle size, the characteristic diameters $d_{60}$, $d_{50}$, $d_{20}$, and $d_{10}$, the uniformity coefficient $U_c = d_{60}/d_{10}$, the specific gravity of soil particles $\rho_s$, the natural water content, and a soil classification. A point of central importance to the prior derivation is that the representative diameter $d_{70}$, which the revised Sellmeijer model \parencite{sellmeijer_2011} requires for the scale factor $F_s$, is not reported directly anywhere in the dataset; its treatment, which is a question of model applicability rather than of routine gap-filling, is addressed in Section~\ref{subsec: Model Applicability at Gravel Grain Sizes}. The consistency tests are blank for the gravelly samples, as expected for cohesionless material. The permeability tests yield a saturated conductivity from reconstituted, density-adjusted specimens, and the triaxial tests yield total-stress and (bracketed) effective-stress cohesion and friction angle, together with the void ratio and bulk density used below to derive the submerged unit weight. In addition, in-situ permeability tests were performed at one depth per borehole, using injection, recovery, or pump-out methods as dictated by local conditions.

\subsection{Grain-Size Results for the Aquifer}
\label{subsec: Grain-Size Results for the Aquifer}

Representative grain-size results for the $A_g$ aquifer samples are reproduced in Table~\ref{tab:grainsize}. The aquifer is consistently classified as a sand-gravel of high uniformity coefficient, with framework $d_{60}$ values ranging from below 1~mm in the fine sandy upper $A_g$ at KP~60.0 to in excess of 10~mm in the coarse gravels at KP~58.8 and KP~62.0. The shallow, sand-dominated samples at the blanket-aquifer transition resolve the finer matrix fraction separately from this coarse framework. This bimodal grading is the empirical basis both for the aquifer conductivity priors and, critically, for the choice of operative representative diameter $d_{70}$ examined in Section~\ref{subsec: Model Applicability at Gravel Grain Sizes}.

\begin{table}[htbp]
\centering
\footnotesize
\caption{Representative grain-size results for $A_g$ aquifer and shallow transition-zone samples \parencite{oyo_1999}. Composition is gravel/sand/silt/clay (per cent). Laboratory $k_s$ is the density-adjusted constant-head value; a dash denotes no permeability test on that specimen. The shallow sand-dominated samples (B-2-1, B-4-1, B-6-1) represent the blanket-aquifer transition and are the basis for the fine-matrix $d_{70}$ interpretation of Section~\ref{subsec: Model Applicability at Gravel Grain Sizes}.}
\label{tab:grainsize}
\resizebox{\textwidth}{!}{%
\begin{tabular}{llccccccc}
\toprule
KP & Sample & Depth (m) & G/S/M/C (\%) & $d_{60}$ (mm) & $d_{50}$ (mm) & $d_{10}$ (mm) & $U_c$ & Lab $k_s$ (m/s) \\
\midrule
57.4 & B-2-1\textsuperscript{\dag} & 0.50--1.00 & 30.4/51.3/13.0/5.3 & 0.635 & 0.365 & 0.0190 & 33.4 & --                  \\
57.4 & B-2-2 & 1.70--4.00 & 38.0/30.0/22.8/9.2 & 1.47  & 0.357 & 0.0057 & 258  & $2.69\times10^{-5}$ \\
57.4 & B-2-3 & 5.00--6.00 & 53.8/35.5/6.5/4.2  & 4.76  & 2.47  & 0.0517 & 92.1 & --                  \\
58.8 & B-4-1\textsuperscript{\dag} & 0.50--1.00 & 24.0/49.7/18.1/8.2 & 0.459 & 0.301 & 0.0071 & 64.6 & --                  \\
58.8 & B-4-2 & 2.30--3.00 & 53.1/28.8/13.1/5.0 & 7.10  & 2.97  & 0.0152 & 467  & $9.16\times10^{-6}$ \\
58.8 & B-4-3 & 6.50--7.00 & 68.9/24.0/4.8/2.3  & 11.1  & 7.62  & 0.192  & 57.8 & $2.84\times10^{-5}$ \\
60.0 & B-6-1\textsuperscript{\dag} & 0.50--1.00 & 4.2/59.6/26.6/9.6  & 0.228 & 0.154 & 0.0053 & 43.0 & $5.59\times10^{-4}$ \\
60.0 & B-6-2 & 2.00--3.00 & 85.3/9.8/3.4/1.5   & 13.2  & 10.5  & 0.530  & 24.9 & $1.83\times10^{-4}$ \\
62.0 & B-9-1 & 1.00--2.00 & 43.9/37.5/13.7/4.9 & 3.46  & 0.906 & 0.0160 & 216  & $1.08\times10^{-5}$ \\
62.0 & B-9-2 & 5.00--6.00 & 73.2/20.5/4.2/2.1  & 11.7  & 7.16  & 0.228  & 51.3 & --                  \\
63.4 & B-2-2 & 1.50--2.00 & 57.0/40.0/3.0/--   & 10.92 & 5.36  & 0.240  & 46.0 & $7.14\times10^{-6}$ \\
63.4 & B-3-2\textsuperscript{\dag} & 1.50--2.00 & 39.0/54.0/7.0/--   & 1.67  & 0.63  & 0.110  & 14.9 & $1.96\times10^{-5}$ \\
\bottomrule
\end{tabular}%
}

\vspace{0.5em}
{\footnotesize $^\dag$~Shallow blanket-aquifer transition sample (gravel fraction $\leq 4$--31\%; sand-dominated). These specimens resolve the fine sand matrix and provide the empirical basis for the operative $d_{70}$ prior (Section~\ref{subsec: Model Applicability at Gravel Grain Sizes}); the deeper samples at each section resolve the coarse gravel framework.}
\end{table}

\subsection{Field Permeability and the Field-to-Laboratory Discrepancy}
\label{subsec: Field Permeability and the Field-to-Laboratory Discrepancy}

The in-situ permeability results for the $A_g$ aquifer are reproduced in Table~\ref{tab:fieldperm}. A systematic and important pattern is that the field conductivity exceeds the density-adjusted laboratory value at the same borehole by factors of approximately 2 to 50. This discrepancy reflects the destruction of in-situ structure and macro-pore connectivity in the reconstituted laboratory specimens, with the consequence that the laboratory values systematically underestimate the bulk aquifer conductivity that governs both the leakage length and the pipe progression rate. This finding is the principal reason that the prior aquifer-conductivity means in Section~\ref{sec: Prior Parameter Distributions: Derivation and Statistical Fitting} are anchored not to the laboratory tests but to the synthesised analysis constants of Section~\ref{subsec: Analysis Constants and the Official Ratings}. The single recovery-method value at KP~63.4 of $4.2\times10^{-3}$~m/s is regarded as a likely outlier and is not used to anchor the prior.

\begin{table}[htbp]
\centering
\footnotesize
\caption{Field (in-situ) permeability of the $A_g$ aquifer \parencite{oyo_1999}. Values consistently exceed the density-adjusted laboratory $k_s$ at the same borehole by factors of roughly 2 to 50.}
\label{tab:fieldperm}
\begin{tabular}{llcll}
\toprule
KP & Borehole & Depth (m) & Method & Field $k_s$ (m/s) \\
\midrule
57.4 & B-2 & 9.00--9.50 & injection  & $5.1\times10^{-5}$ \\
58.8 & B-4 & 8.00--8.50 & injection  & $2.2\times10^{-6}$ \\
60.0 & B-6 & 8.40--8.90 & pump-out   & $1.2\times10^{-4}$ \\
62.0 & B-9 & 7.40--7.90 & pump-out   & $7.0\times10^{-5}$ \\
63.4 & B-2 & $\approx5.5$ & recovery & $4.2\times10^{-3}$ (outlier) \\
63.4 & B-3 & $\approx5.7$ & recovery & $6.3\times10^{-5}$ \\
\bottomrule
\end{tabular}
\end{table}

A mineralogical caveat applies to the use of these framework-derived specific gravities as the empirical basis for the Sellmeijer resistance factor $F_r$. The specific gravity values of 2.64 to 2.77 are consistent with dense granitic or quartzitic coarse clasts, which dominate the gravel framework of the $A_g$ specimens tested. The Sellmeijer formulation, however, treats $\gamma'_s$ as the submerged weight of the eroding grain at the pipe tip, which under the matrix-controlled interpretation of Section~\ref{subsec: Model Applicability at Gravel Grain Sizes} is the sand matrix fraction rather than the gravel framework. The Tokachi Plain sediments are documented as poorly sorted periglacial and fluvial deposits derived from the Hidaka granitic and volcanic ranges \parencite{furuichi_2018}, and the matrix fines may contain a proportion of pumiceous or vesicular volcanic-glass particles whose effective specific gravity can fall substantially below the quartz-sand reference of approximately 2.65 that underlies the Sellmeijer calibration dataset. A pumiceous grain with $\rho_s$ in the range 1.8 to 2.3~t/m\textsuperscript{3} would reduce the buoyant grain weight entering $F_r$ by roughly 25 to 55 per cent relative to the quartz basis, lowering the predicted critical head $H_c$ by a similar proportion and increasing the conditional failure probability accordingly. The present $\gamma'_s$ prior, with means in the range 11.7 to 12.7~kN/m\textsuperscript{3} tabulated in Section~\ref{subsec: Stochastic Parameter Vector}, is derived entirely from the coarse gravel triaxials and may therefore overstate the matrix grain weight that physically governs erosion resistance in the Pol--Sellmeijer framework. This limitation cannot be resolved without matrix-fraction petrographic analysis for the governing sections KP~62.0 and KP~63.4, where the blanket is thinnest and $F_r$ is most influential on the critical head. Pending such analysis, the uncertainty is bounded by evaluating Phase~1 fragility curves under both the measured gravel-framework $\gamma'_s$ and a lower-bound matrix estimate consistent with a 20 per cent pumiceous admixture, a contrast that is carried as a sensitivity result in Chapter~\ref{chap: Results: Subsurface Piping Assessment} and addressed in Layer~2 of the Discussion.

\subsection{Analysis Constants and the Official Ratings}
\label{subsec: Analysis Constants and the Official Ratings}

A distinctive feature of the OYO report is that it records, separately from the raw test data, the set of specified soil constants that were actually used in the seepage and slope computations and that produced the official 1998 ratings of Table~\ref{tab:oyo_1998}. These analysis constants, extracted for the blanket and aquifer units, are summarised in Table~\ref{tab:form5}. They are central to the present work because they represent the project engineers' considered synthesis of the multi-method permeability evidence, reconciling the laboratory tests, the field tests, and the standard penetration record into a single value per layer. As such they sit at the high end of the conductivity evidence, consistent with the field-to-laboratory discrepancy of Section~\ref{subsec: Field Permeability and the Field-to-Laboratory Discrepancy}, and they are adopted as the means of the conductivity priors in Section~\ref{sec: Prior Parameter Distributions: Derivation and Statistical Fitting}.

\begin{table}[htbp]
\centering
\footnotesize
\caption{Specified soil constants for the blanket ($A_c$) and aquifer ($A_g$) units used in the 1998 seepage computations \parencite{oyo_1999}. Conductivities converted from the recorded values; unit weights from the recorded saturated densities with $\gamma_w = 9.81$~kN/m\textsuperscript{3}. At KP~63.4 the blanket is collapsed into a single foundation unit and no separate $A_c$ value exists.}
\label{tab:form5}
\begin{tabular}{lcccc}
\toprule
KP & $A_c$ blanket $k_{bl}$ (m/s) & $A_g$ aquifer $k_\mathrm{aq}$ (m/s) & $A_c$ sat. density (t/m\textsuperscript{3}) & $A_g$ sat. density (t/m\textsuperscript{3}) \\
\midrule
57.40 & $1.6\times10^{-6}$ & $3.0\times10^{-3}$ & 1.70 & 1.98 \\
58.80 & $1.0\times10^{-6}$ & $2.0\times10^{-3}$ & 1.70 & 2.02 \\
60.00 & $1.0\times10^{-6}$ & $1.0\times10^{-3}$ & 1.70 & 2.07 \\
62.00 & $3.0\times10^{-6}$ & $1.0\times10^{-3}$ & 1.70 & 2.18 \\
63.40 & n/a                & $6.0\times10^{-5}$ & n/a  & 2.00 \\
\bottomrule
\end{tabular}
\end{table}

\subsection{Design Hydrograph Parameters}
\label{subsec: Geotechnical Dataset Design Hydrograph}

The OYO report also specifies the external loading conditions used in the 1998 seepage computations. The rainfall forcing is common to all sections, comprising an antecedent rainfall of 100~mm over 100~hours followed by a flood-period rainfall of 246~mm over 24.6~hours. The resulting design water-level hydrographs are referenced to the Obihiro gauge for KP~57.4, KP~58.8, KP~60.0, and KP~62.0, all sharing an identical waveform with a peak plateau of 1.5~hours and a recession rate of 0.29~m/h, and to the Shikaribetsu River confluence gauge for KP~63.4, which uses a shorter event with no plateau and a recession rate of 0.20~m/h. The Obihiro-referenced common waveform is confirmed by the KP~57.4 water-level time series, which holds the high-water level for 1.5~hours before receding 5.51~metres over 19~hours, equal to 0.29~m/h. This shared loading across four of the five sections is what permits the section-to-section contrasts discussed in Section~\ref{subsec: Cross-Section Inventory and the Foreshore-Width Control on Risk} to be attributed to geometry and stratigraphy rather than to differences in forcing. It must be emphasised, however, that these OYO design events are reference checks against the historical deterministic framework only; the operative loading throughout this thesis is the d4PDF climate ensemble described in Section~\ref{sec: The d4PDF Climate Ensemble: Structure, Downscaling, and Hydrograph Extraction}, which carries its own epistemic limitation in that it was processed on an annual-maximum framework calibrated for peak discharge rather than for the load duration that governs BEP.

\section{Prior Parameter Distributions: Derivation and Statistical Fitting}
\label{sec: Prior Parameter Distributions: Derivation and Statistical Fitting}

\subsection{Distributional Family and Lognormal Transform}
\label{subsec: Distributional Family and Lognormal Transform}

The stochastic geotechnical parameters are represented by lognormal distributions, consistent with the convention of the reliability framework of \textcite{pol_sie_2024} and with standard practice in probabilistic levee assessment. For a parameter $X$ with arithmetic mean $\mu$ and coefficient of variation $\mathrm{CoV}$, the parameters of the underlying normal distribution of $\ln X$ are
\begin{equation}
\sigma_{\ln} = \sqrt{\ln\!\left(1 + \mathrm{CoV}^2\right)},
\qquad
\mu_{\ln} = \ln(\mu) - \tfrac{1}{2}\sigma_{\ln}^2,
\end{equation}
which preserves the arithmetic mean under back-transformation, $\mathbb{E}[X] = \exp(\mu_{\ln} + \tfrac{1}{2}\sigma_{\ln}^2) = \mu$. All location parameters $\mu_{\ln}$ are evaluated in SI base units (metres, metres per second, kilonewtons per cubic metre) to ensure dimensional consistency in the Monte Carlo engine. As a worked illustration, the aquifer conductivity at KP~57.4, with $\mu = 3.0\times10^{-3}$~m/s and $\mathrm{CoV} = 0.50$, gives $\sigma_{\ln} = \sqrt{\ln 1.25} = 0.472$ and $\mu_{\ln} = \ln(3.0\times10^{-3}) - \tfrac{1}{2}(0.472)^2 = -5.921$, the value tabulated in Table~\ref{tab:priors_muln}.

Throughout this section, the prior means are sourced from the OYO dataset as detailed in Sections~\ref{sec: Geotechnical Dataset: Borehole Logs and Laboratory Testing}, while the coefficients of variation are adopted, wherever a direct correspondence exists, from the base-case reliability parameterization of \textcite{pol_sie_2024}. Where no direct correspondence exists, the adopted value is justified by analogue and flagged explicitly in Section~\ref{subsec: Treatment of Data Gaps and Judgment Calls}.

\subsection{The Stochastic Parameter Vector}
\label{subsec: Stochastic Parameter Vector}

The stochastic parameter vector comprises the seven random variables
\begin{equation}
\boldsymbol{\theta} =
\bigl(k_\mathrm{aquifer},\; d_{70},\; D_\mathrm{aq},\; D_{bl},\;
      k_{bl},\; \gamma'_\mathrm{bl},\; C_e\bigr),
\end{equation}
of which the first six are time-invariant geotechnical quantities and the seventh is the empirical erosion coefficient of the Pol progression ODE. The distributional specification, common to all cross-sections, is given in Table~\ref{tab:priors_phase1}; the KP-specific means and the resulting lognormal location parameters are given side by side in Tables~\ref{tab:priors_means} and~\ref{tab:priors_muln}.

\begin{table}[htbp]
\centering
\scriptsize
\renewcommand{\arraystretch}{0.85}
\caption{Distributional specification of the stochastic parameter vector $\boldsymbol{\theta}$ (common to all cross-sections). Means are sourced from the OYO dataset as indicated; CoV values are adopted from \textcite{pol_sie_2024} where a direct correspondence exists, and by analogue otherwise (see Section~\ref{subsec: Treatment of Data Gaps and Judgment Calls}).}
\label{tab:priors_phase1}
\resizebox{\textwidth}{!}{%
\begin{tabular}{llccll}
\toprule
Parameter & Family & CoV & $\sigma_{\ln}$ & Mean source & CoV source \\
\midrule
$k_\mathrm{aquifer}$ & Lognormal & 0.50  & 0.472 & OYO analysis constant ($A_g$)     & \textcite{pol_sie_2024} \\
$d_{70}$             & Lognormal & 0.30  & 0.294 & OYO $d_{60}$ (sand matrix)         & widened (see text) \\
$D_\mathrm{aq}$      & Lognormal & 0.10  & 0.100 & OYO borehole panels ($A_g$/$N_s$) & \textcite{pol_sie_2024}, rescaled \\
$D_{bl}$             & Lognormal & 0.167 & 0.166 & OYO borehole panels ($A_c$)       & \textcite{pol_sie_2024}, absolute $\sigma$ \\
$k_{bl}$             & Lognormal & 0.50  & 0.472 & OYO analysis constant ($A_c$)     & analogue to $k_\mathrm{aquifer}$ \\
$\gamma'_\mathrm{bl}$ & Lognormal & 0.056 & 0.056 & OYO Form~5 $\rho_t$ ($A_c$), $\gamma'_\mathrm{bl} = \gamma_{sat,bl} - \gamma_w$ & \textcite{pol_sie_2024} ($A_c$ blanket parameterization) \\
$C_e$                & Lognormal & 0.78  & 0.691 & Pol SIE 2024 base case (0.055) & \textcite{pol_sie_2024}, Tab.~2 \\
\bottomrule
\end{tabular}%
}
\end{table}

\noindent
\resizebox{\textwidth}{!}{%
  \begin{tabular}{@{} c @{\hspace{1em}} c @{}}

    \begin{minipage}[b]{8.2cm}
      \captionof{table}{Cross-section-specific prior means for $\boldsymbol{\theta}$. Aquifer and blanket conductivities from OYO analysis constants (Table~\ref{tab:form5}); thicknesses from Table~\ref{tab:strat_thickness}; $d_{70}$ is the fine sand-matrix diameter (Section~\ref{subsec: Model Applicability at Gravel Grain Sizes}), bulk-gravel values in brackets as co-primary result; $\gamma'_\mathrm{bl}$ from the OYO Form~5 $A_c$ design density ($\gamma'_\mathrm{bl} = \gamma_{sat,bl} - \gamma_w$), basin-wide for KP~57.4--62.0; n/a at KP~63.4 ($A_c$ absent); $C_e$ basin-wide.}
      \label{tab:priors_means}
    \end{minipage}
    &
    \begin{minipage}[b]{8.2cm}
      \captionof{table}{Resulting lognormal location parameters $\mu_{\ln}$ for $\boldsymbol{\theta}$, evaluated in SI base units and paired with the section-invariant $\sigma_{\ln}$ of Table~\ref{tab:priors_phase1}.}
      \label{tab:priors_muln}
    \end{minipage} \\[-1.0em]

    {\scriptsize
     \renewcommand{\arraystretch}{0.85}
     \setlength{\tabcolsep}{3.5pt}
     \begin{tabular}[t]{lccccc}
      \toprule
      Parameter (unit) & KP~57.4 & KP~58.8 & KP~60.0 & KP~62.0 & KP~63.4 \\
      \midrule
      $k_\mathrm{aquifer}$ (m/s) & $3.0\times10^{-3}$ & $2.0\times10^{-3}$ & $1.0\times10^{-3}$ & $1.0\times10^{-3}$ & $6.0\times10^{-5}$ \\
      $d_{70}$ (mm)              & 0.70  & 0.53  & 0.26  & 0.70  & 0.70  \\
      \quad{\tiny[bulk, co-primary]}  & {\tiny 5.5} & {\tiny 13} & {\tiny 1.3} & {\tiny 13.5} & {\tiny 9.5} \\
      $D_\mathrm{aq}$ (m)        & 7     & 8     & 9     & 10    & 11    \\
      $D_{bl}$ (m)               & 0.80  & 0.85  & 0.85  & 0.45  & 1.0   \\
      $k_{bl}$ (m/s)             & $1.6\times10^{-6}$ & $1.0\times10^{-6}$ & $1.0\times10^{-6}$ & $3.0\times10^{-6}$ & n/a \\
      $\gamma'_\mathrm{bl}$ (kN/m$^3$) & 6.9 & 6.9 & 6.9 & 6.9 & n/a \\
      $C_e$ (--)                 & \multicolumn{5}{c}{0.055 (basin-wide)} \\
      \bottomrule
      \end{tabular}%
    }
    &
    {\scriptsize
     \renewcommand{\arraystretch}{0.85}
     \setlength{\tabcolsep}{3.5pt}
     \begin{tabular}[t]{lccccc}
      \toprule
      Parameter & KP~57.4 & KP~58.8 & KP~60.0 & KP~62.0 & KP~63.4 \\
      \midrule
      $k_\mathrm{aquifer}$ & $-5.921$  & $-6.326$  & $-7.019$  & $-7.019$  & $-9.833$ \\
      $d_{70}$             & $-7.308$  & $-7.586$  & $-8.298$  & $-7.308$  & $-7.308$ \\
      $D_\mathrm{aq}$      & $1.941$   & $2.074$   & $2.192$   & $2.298$   & $2.393$  \\
      $D_{bl}$             & $-0.237$  & $-0.176$  & $-0.176$  & $-0.812$  & $-0.014$ \\
      $k_{bl}$             & $-13.457$ & $-13.927$ & $-13.927$ & $-12.829$ & n/a      \\
      $\gamma'_\mathrm{bl}$ & $1.930$   & $1.930$   & $1.930$   & $1.930$   & n/a      \\
      $C_e$                & \multicolumn{5}{c}{$-3.139$ (basin-wide)} \\
      \bottomrule
      \end{tabular}%
    }
  \end{tabular}%
}

The physical role of each variable in the limit states of Chapter~\ref{chap: Methodology} motivates its treatment as stochastic. The aquifer conductivity $k_\mathrm{aquifer}$ governs both the leakage length and, through its appearance inside the overloading term of the Pol progression ODE, the pipe progression velocity \parencite{pol_compgeo_2024}; its mean is taken from the synthesised analysis constant of Table~\ref{tab:form5} rather than from the laboratory tests, for the reasons established in Section~\ref{subsec: Field Permeability and the Field-to-Laboratory Discrepancy}. The representative diameter $d_{70}$ governs grain entrainment resistance through the Sellmeijer scale factor $F_s$ and the equilibrium head; because the dataset reports no $d_{70}$ and the choice of representative fraction in a sand-gravel is decisive for the predicted critical head, its derivation is treated separately and at length in Section~\ref{subsec: Model Applicability at Gravel Grain Sizes}, where the fine sand-matrix diameter is adopted as the operative prior mean, with the bulk-gravel value retained as a co-primary result for reasons developed in that section and in Section~\ref{subsubsec: k-d70 cross-correlation}. The aquifer thickness $D_\mathrm{aq}$ governs both the leakage length and the critical pipe length $l_c$, and is read directly from the borehole panels. The blanket thickness $D_{bl}$ and conductivity $k_{bl}$, together with the submerged bulk unit weight of the confining blanket $\gamma'_\mathrm{bl}$, jointly govern the leakage length and the uplift and heave limit states. The latter is derived from the $A_c$ design density of the OYO analysis constants (Table~\ref{tab:form5}), which gives the saturated bulk unit weight $\gamma_{sat,bl} = \rho_t\, g = 1.70 \times 9.81 = 16.7$~kN/m\textsuperscript{3}, uniformly across KP~57.4--62.0. The submerged bulk unit weight follows by subtraction of $\gamma_w$:
\begin{equation}
\gamma'_\mathrm{bl} = \gamma_{sat,bl} - \gamma_w = 6.9~\mathrm{kN/m^3},
\end{equation}
with a coefficient of variation of $0.056$ adopted directly from the blanket parameterization of \textcite{pol_sie_2024}. At KP~63.4 the $A_c$ blanket is absent (Section~\ref{sec: app form5}) and $\gamma'_\mathrm{bl}$ is correspondingly undefined for that section, consistent with the treatment of $k_{bl}$. It is essential to note that $\gamma'_\mathrm{bl}$ is the submerged weight of the confining blanket material, entering only the uplift limit state $Z_\mathrm{uplift}$ and the critical heave gradient $i_c$ in $Z_\mathrm{heave}$ (Chapter~\ref{chap: Methodology}); it is distinct from the submerged unit weight of the aquifer sand particles, $\gamma'_\mathrm{p}$, which enters the Sellmeijer resistance factor $F_r$, the critical head $H_c$, and the equilibrium head $H_\mathrm{eq}$, and which is held deterministic at the basin-wide value derived in Section~\ref{subsec: Fixed Parameters Prior}. The erosion coefficient $C_e$ is treated as basin-wide and is discussed below.

The seepage length $L$ is carried as an additional stochastic geometric parameter, sampled per cross-section independently of $\boldsymbol{\theta}$ under a per-section lognormal prior. The angle of repose, the relative density, the uniformity coefficient, and the grain angularity are held deterministic, the latter three being set to their experimental mean values so that their contributions to the Sellmeijer resistance factor reduce to unity, following the conventions of \textcite{sellmeijer_2011} and \textcite{pol_sie_2024}. The exit elevation is the surveyed landside toe elevation $z_\mathrm{toe}$, at which the confined aquifer overpressure is released at the natural ground surface; the earlier riverside foreshore-crest placeholder is superseded.

\subsection{Fixed Parameters}
\label{subsec: Fixed Parameters Prior}

Three parameters of the Sellmeijer grain-stability formulation are held deterministic at their Pol base-case or site-derived values: the angle of repose $\theta = 37^\circ$, White's coefficient $\eta = 0.25$ \parencite{pol_sie_2024}, and the submerged unit weight of the aquifer sand particles $\gamma'_\mathrm{p}$. The latter is computed as $\gamma'_\mathrm{p} = (G_s-1)\gamma_w$ from the specific gravities of the $A_g$ specimens in Table~\ref{tab:grainsize} ($n=12$, KP~57.4--63.4), giving a sample mean of $16.87$~kN/m\textsuperscript{3} with a coefficient of variation of approximately $1.6\%$. This near-zero variability, comparable to that of $D_r$, $C_{u,m}$, and $\mathrm{KAS}_m$ below, justifies treating $\gamma'_\mathrm{p}$ as a deterministic constant rather than as an eighth stochastic dimension. Note that $\gamma'_\mathrm{p}$ (particle-scale, no porosity correction) is distinct from the blanket bulk unit weight $\gamma'_\mathrm{bl}$ of Table~\ref{tab:priors_phase1}, which retains its place in the stochastic vector because it governs the uplift and heave limit states rather than the Sellmeijer resistance factors.  In addition, and as noted above, the relative density $D_r = 0.725$, the uniformity coefficient $C_{u,m} = 1.81$, and the angularity $\mathrm{KAS}_m = 0.498$ are fixed at the experimental means of the multivariate regression of \textcite{sellmeijer_2011}, so that the corresponding ratio terms in $F_r$ equal unity. This treatment isolates the active geotechnical uncertainty in the six stochastic geotechnical variables of Table~\ref{tab:priors_phase1} and prevents the spurious propagation of variability through parameters for which no section-specific measurement is available, or whose measured variability is negligible.

\subsection{Supplementary Parameters for the Uplift and Heave Limit States}
\label{subsec: Supplementary Parameters for the Uplift and Heave Limit States}

The submerged blanket bulk unit weight $\gamma'\mathrm{bl}$ that governs the uplift and heave limit states is carried directly in the stochastic vector $\boldsymbol{\theta}$ (Table~\ref{tab:priors_phase1}) rather than as a supplementary parameter, resolving the symbol collision that previously existed between the aquifer grain weight and the blanket weight under the shared symbol $\gamma's$. No additional model-uncertainty factors are applied to the uplift or heave limit states. Pol's uplift model factor $m_u$ and critical-head model factor $m_p$ \parencite{pol_sie_2024} are not carried in the baseline parameterization: the single stochastic quantity of the progression law is the erosion coefficient $C_e$, whose uncertainty the Phase~2 survival constraint is able to update, and adding several further uncalibrated factors that the available data cannot separately constrain is deferred to a sensitivity rather than folded into the baseline. This is not a claim that $C_e$ absorbs the model-form uncertainty those factors represent; consistent with the treatment of $C_e$ in Section~\ref{subsec: Stochastic Parameter Vector}, the laminar-versus-turbulent uncertainty of the critical head is nominally $m_p$'s to carry, and $m_p \sim \mathrm{Lognormal}(1, 0.12)$ is available as an explicit opt-in factor whose effect on the fragility is quantified separately as a robustness check rather than presented inside the baseline curves. The critical heave gradient is likewise not carried as a separate calibrated input: consistent with the formulation $i_c = \gamma'\mathrm{bl}/\gamma_w$ of Chapter~\ref{chap: Methodology} \parencite{terzaghi1943}, it is computed within each Monte Carlo realisation directly from the sampled $\gamma'\mathrm{bl}$, so that its uncertainty is inherited entirely from $\gamma'_\mathrm{bl}$'s coefficient of variation of $0.056$ rather than carried as an independently specified $0.143$.

\subsection{Model Applicability at Gravel Grain Sizes: The Operative \texorpdfstring{$d_{70}$}{d70}}
\label{subsec: Model Applicability at Gravel Grain Sizes}

The treatment of the representative grain diameter $d_{70}$ is not a routine gap-filling exercise but a question of model applicability, and it is the single most consequential judgment in the prior derivation. The OYO dataset reports no $d_{70}$ on any specimen; only $d_{60}$, $d_{50}$, $d_{20}$, and $d_{10}$ are tabulated. A naive interpretation extrapolates $d_{70}$ from the reported $d_{60}$ of the representative $A_g$ samples, giving the values shown in the bulk-bound row of Table~\ref{tab:priors_means}, namely 1.3 to 13.5~mm. These values lie one to almost two orders of magnitude above the calibration basis of the revised Sellmeijer model, whose multivariate regression was fitted on sands with an experimental mean diameter of 0.208~mm and a validated parameter range extending only to 0.430~mm \parencite{sellmeijer_2011}. At a bulk diameter of 13~mm the critical-head formula is therefore extrapolating far outside its experimental basis, and the extrapolation is not benign: the scale factor carries a net dependence $F_s \propto d_{70}^{0.4}$, so that $H_c$ scales as approximately $d_{70}^{0.4}$, and the coarse-sand behaviour that \textcite{sellmeijer_2011} themselves flag as poorly understood (their single coarse-sand validation test deviating by 25~per cent) is precisely the regime a bulk-gravel $d_{70}$ would invoke. Adopting the bulk diameter would inflate the predicted critical head by a factor of roughly 1.9 to 3.6 relative to the matrix interpretation developed below, which would in turn suppress the conditional fragility and risk presenting the study sections as artificially immune to piping.

The physical resolution is to recognise that backward erosion in a bimodal sand-gravel is initiated and sustained by the erosion of the finer matrix fraction rather than of the gravel framework. Once the sand matrix is removed by the seepage flow, the coarse clasts are no longer supported and are transported as a lag or raft, so the grain size that governs the entrainment resistance at the pipe tip is the matrix diameter, not the framework diameter. This matrix-controlled interpretation of piping in gravelly aquifers is consistent with the grain-stability basis of the Sellmeijer formulation, in which the eroding particle is the one that must resist the flow forces while the finer particles around it are carried away \parencite{sellmeijer_2011}, and it is the interpretation adopted here. The operative $d_{70}$ is accordingly taken from the fine sand-matrix evidence in the shallow blanket-aquifer transition samples, which are the sand-dominated specimens B-2-1 at KP~57.4, B-4-1 at KP~58.8 ($d_{60} = 0.459$~mm), and B-6-1 at KP~60.0 ($d_{60} = 0.228$~mm). The resulting matrix means of 0.26 to 0.70~mm, tabulated in Table~\ref{tab:priors_means}, lie within or immediately adjacent to the Sellmeijer validated range, with KP~60.0 falling squarely inside it. The coefficient of variation is widened from the Pol base-case value of 0.10 to 0.30, reflecting the genuine within-section grading heterogeneity evident in Table~\ref{tab:grainsize} between the matrix and framework fractions.

The matrix-controlled interpretation resolves the model-applicability problem for $d_{70}$, but it introduces a coupled internal inconsistency in the parameter vector that must be addressed before sampling begins. The Sellmeijer scale factor $F_s$ contains the kinematic permeability coefficient $\kappa = k\nu/g$, so $k_\mathrm{aquifer}$ and $d_{70}$ appear together inside the same physics; the Pol progression ODE likewise carries $k_\mathrm{aquifer}$ to the power 0.81 in the pipe-growth velocity \parencite{pol_compgeo_2024}. A sand matrix with $d_{70} \approx 0.26$--0.70~mm is inconsistent, via the Hazen or Kozeny--Carman relations, with a conductivity of order $10^{-3}$~m/s: the Hazen estimate $k \approx C_H d_{10}^2$ for a 0.3~mm-class matrix gives $k \sim 10^{-4}$~m/s, roughly one order of magnitude below the bulk-gravel analysis constants anchoring $k_\mathrm{aquifer}$. Sampling the two parameters as statistically independent quantities therefore risks constructing realisations in which the matrix grain governs the resistance physics while the gravel framework governs the progression velocity, a pairing that describes two physically incompatible soils within a single Monte Carlo draw. When both parameters simultaneously occupy their upper prior tails under independent sampling, this incoherence inflates the prior transient fragility beyond what either physical interpretation alone would predict: high $k_\mathrm{aquifer}$ drives fast pipe progression through the ODE, while low matrix $d_{70}$ lowers the critical head, so both individually push prior failure upward. A disproportionate share of the apparent Bayesian updating in Phase~2 may then reflect the elimination of physically impossible realisations rather than the genuinely informative conditioning on the 2016 survival event, a distortion that the architecture specification identifies as Failure Mode~7 for the analogous $C_e$--$k_\mathrm{aquifer}$ interaction, and that applies equally, and more fundamentally, here.

Two remediation strategies are available, and the choice between them is determined empirically by Diagnostic~A of the pre-analysis plan (Section~\ref{subsubsec: k-d70 cross-correlation}). The first, and preferred, strategy is to impose an explicit cross-correlation $\rho(\ln k_\mathrm{aquifer},\, \ln d_{70})$ through a Nataf (Gaussian copula) transform in the prior sampler, with $\rho$ estimated by regressing $\ln k_\mathrm{lab}$ against $\ln d_{10}$ (or $\ln d_{60}$) on the paired permeability and grain-size specimens available in Tables~\ref{tab:grainsize} and~\ref{tab:form5}. A statistically meaningful, Hazen-consistent slope from this regression would confirm that conductivity and grain size co-vary within a single population and that a Nataf correlation is the appropriate link. The second strategy applies if the regression reveals a flat or bimodal relationship, indicating that the framework conductivities and the matrix grain sizes belong to distinct, physically decoupled populations, in which case the two-soil model is adopted explicitly: $d_{70}$ (matrix fraction) enters the Sellmeijer resistance factors and equilibrium head, while $k_\mathrm{aquifer}$ (framework) governs the seepage rate and progression velocity, and the two parameters are drawn from separate marginals without a Nataf link. The architecture specification, Section~7, is built to receive a full $7 \times 7$ correlation matrix $\boldsymbol{\Sigma}$ as a configuration input, so either strategy is accommodated without structural change to the engine; the empirical $\rho$, or the two-soil flag, is supplied once Diagnostic~A is complete. In either case, independent sampling of $k_\mathrm{aquifer}$ and $d_{70}$ as currently implied by the identity default in Section~7 is not a permissible outcome of this diagnostic: the coherence question must be resolved, and its resolution must be recorded as a configuration parameter before Phase~1 runs.

The bulk-gravel diameters are therefore retained, in the bracketed row of Table~\ref{tab:priors_means}, as co-primary results of Phase~1 alongside the matrix interpretation, not as a secondary sensitivity bound. Reporting both is essential because the choice of operative $d_{70}$ propagates directly into the critical head and thence into the prior fragility curves; the contrast between the two interpretations is itself a quantitative statement of how far the Sellmeijer model can be trusted in this gravel-dominated setting, and the decomposition of this uncertainty is a required input to the Layer~4 discussion of the prior-to-posterior shift in Chapter~\ref{chap: Discussion}.

\subsection{Treatment of Data Gaps and Judgment Calls}
\label{subsec: Treatment of Data Gaps and Judgment Calls}

The prior distributions tabulated above rest on a set of further interpretive decisions that must be made explicit, both for transparency and because several of them directly condition the calibrated posterior of Phase~2. The treatment of $d_{70}$, being a question of model applicability rather than of parameter uncertainty alone, has been addressed separately in Section~\ref{subsec: Model Applicability at Gravel Grain Sizes}.

\subsubsection*{Conductivity priors and the field-to-laboratory discrepancy.}
The aquifer and blanket conductivity means are anchored to the synthesised analysis constants of Table~\ref{tab:form5} rather than to the laboratory tests, because, as established in Section~\ref{subsec: Field Permeability and the Field-to-Laboratory Discrepancy}, the laboratory values underestimate the bulk conductivity by factors of roughly 2 to 50. These analysis constants sit at the high end of the conductivity evidence, which is conservative for transient BEP in isolation, since a higher $k_\mathrm{aquifer}$ drives faster pipe progression through the Pol ODE. This reasoning is, however, qualified by the cross-correlation consideration developed in Section~\ref{subsec: Model Applicability at Gravel Grain Sizes}: when $k_\mathrm{aquifer}$ is sampled jointly with the matrix $d_{70}$ without an imposed correlation, high $k_\mathrm{aquifer}$ and low $d_{70}$ can co-occur in the same Monte Carlo realisation in a way that simultaneously inflates progression velocity and lowers the entrainment threshold, describing two incompatible soils in a single draw and overstating the prior transient fragility in a manner that is not simply conservative. The Nataf cross-correlation treatment described in Section~\ref{subsubsec: k-d70 cross-correlation} below mitigates this effect by constraining sampling to the physically coherent joint region of the two marginals. Subject to that constraint, the deliberate anchoring of $k_\mathrm{aquifer}$ to the high-end analysis constants remains internally consistent with the Phase~2 design: it is precisely the tail of high-conductivity, fast-progressing realisations, conditioned on physically plausible $d_{70}$, that the Bayesian Accept-Reject filter is designed to trim against the documented 2016 survival, so that the high-end prior and the survival-conditioned posterior are internally consistent components of the same updating procedure.

\subsubsection*{Cross-correlation between $k_\mathrm{aquifer}$ and $d_{70}$.}
\label{subsubsec: k-d70 cross-correlation}

The physical coupling between aquifer conductivity and representative grain size, established in Section~\ref{subsec: Model Applicability at Gravel Grain Sizes}, requires that $k_\mathrm{aquifer}$ and $d_{70}$ not be sampled as independent random variables. This subsubsection states the pre-analysis diagnostic that resolves the form of their dependence structure, and the two architectural consequences that follow from each possible outcome.

\textit{Diagnostic~A.} The correlation coefficient $\rho(\ln k_\mathrm{aquifer},\, \ln d_{70})$ is estimated by regressing $\ln k_\mathrm{lab}$ against $\ln d_{10}$ and separately against $\ln d_{60}$ on all specimens in Tables~\ref{tab:grainsize} and~\ref{tab:form5} for which both a grain-size descriptor and a laboratory conductivity are concurrently available. The regression yields two diagnostic quantities: (i) the slope, compared against the Hazen prediction of 2 as a physical consistency check, and (ii) the Pearson correlation coefficient $r$, which serves as the candidate Nataf input. A slope within a factor of approximately two of the Hazen prediction, combined with a coefficient of determination $r^2 \gtrsim 0.5$, is taken as evidence that conductivity and grain size co-vary within a single, physically coherent soil population.

\textit{Outcome~1 (Nataf copula).} If Diagnostic~A recovers a meaningful Hazen-consistent slope, the parameters are treated as belonging to a single population. A Nataf (Gaussian copula) transform is imposed in module~M2 of the prior sampler with the empirical correlation coefficient $\rho = r$, and the full $7 \times 7$ correlation matrix $\boldsymbol{\Sigma}$ supplied to the architecture specification Section~7 has $\Sigma_{k,d} = \rho$ and all other off-diagonal entries equal to zero (the remaining parameters are sampled independently). Both the matrix and bulk $d_{70}$ interpretations are then run with the same Nataf structure; the contrast between their resulting fragility curves constitutes the primary quantification of the two-soil uncertainty and is reported as a main result of Phase~1, not a post-hoc sensitivity.

\textit{Outcome~2 (two-soil model).} If the regression is flat ($r^2 < 0.3$) or reveals a bimodal structure, indicating that the laboratory conductivities belong to the gravel framework while the grain sizes characterise the sand matrix, and that the two populations do not co-vary, the two-soil model is adopted. Under this model, $d_{70}$ (matrix fraction) governs the Sellmeijer entrainment-resistance factors $F_r$ and $F_s$ and the equilibrium head $H_\mathrm{eq}$, while $k_\mathrm{aquifer}$ (framework) governs the seepage rate and pipe-growth velocity in the Pol ODE. The two parameters are drawn from their respective marginal lognormal distributions without a Nataf link ($\Sigma_{k,d} = 0$), but their physical roles in separate model components are documented explicitly, and the decoupling is carried as a first-order epistemic limitation in Chapter~\ref{chap: Discussion}, Layer~4.

In both outcomes, the identity default currently implied by the architecture specification Section~7, equivalent to $\boldsymbol{\Sigma} = \mathbf{I}$, that is, fully independent sampling, is superseded. The empirical $\rho$ (Outcome~1) or the two-soil flag (Outcome~2) must be entered as a named configuration parameter before any Phase~1 run is executed. The engine structure accommodates both outcomes without modification, since module~M2 is built to accept an arbitrary $\boldsymbol{\Sigma}$; the analytical work of Diagnostic~A is the sole remaining prerequisite.

\subsubsection*{Coefficients of variation without a direct source.}
Two coefficients of variation are adopted by analogue and flagged accordingly. The blanket conductivity $k_{bl}$ has no entry in the reliability parameterization of \textcite{pol_sie_2024}, which fixes the response factor $r_e$ deterministically; because the present framework instead derives $r_e$ from $k_{bl}$ through the Mazure leakage-length relation, a conductivity coefficient of variation equal to that of $k_\mathrm{aquifer}$ is adopted. The aquifer grain weight $\gamma'_s$ borrows the coefficient of variation of the saturated blanket unit weight from the same source, a choice supported by the empirical spread of the $A_g$ specific gravity and void ratio, which brackets the adopted value of 0.056. For the aquifer thickness $D_\mathrm{aq}$, the Pol base case specifies an absolute standard deviation of 0.5~m on a 20~m reference aquifer, equivalent to a coefficient of variation of 0.025; transferred directly, this would impose an implausibly tight prior on the 7 to 11~m Tokachi aquifers read from scanned borehole panels. A coefficient of variation of 0.10 is therefore adopted, corresponding to an absolute standard deviation of roughly 0.7 to 1.1~m across the sections, which is a more defensible reflection of the uncertainty in picking the aquifer base from the panels and which is preferred over carrying a known-inadequate value into the prior.

\subsubsection*{Erosion coefficient $C_e$.}
The erosion coefficient is adopted as basin-wide with a mean of 0.055 and a standard deviation of 0.043, giving a coefficient of variation of 0.78 and $\mu_{\ln} = -3.139$. This is the field-reliability base case of \textcite{pol_sie_2024} (Table~2), recommended by the model's author for levee reliability assessment as resting on the largest set of validation experiments and as the conservative (higher-failure-probability) choice. It is adopted here in place of the smaller values obtained by direct experimental calibration, for a reason that is itself informative: the two families of values correspond to two \emph{different calibration targets}. Reproducing the detailed \emph{time-dependent} pipe-length development in the individual small- and large-scale experiments requires $C_e \approx 0.014$--$0.016$, with a between-test scatter of approximately 0.007 to 0.030 \parencite{pol_compgeo_2024}; matching instead the \emph{mean post-critical progression rate} across the full set of validation tests, the target underlying the reliability parameterization, requires $C_e \approx 0.055$. The factor of roughly three to four between the two targets is unexplained even by the model author, and the width of the adopted prior (a coefficient of variation of 0.78, comparable to the between-test scatter of the calibrations themselves rather than an order of magnitude tighter) is retained precisely because it reflects this genuine epistemic uncertainty in the coefficient, which the Phase~2 survival constraint is then able to reduce.

Two points of framing follow. First, treating $C_e$ as stochastic is justified by the intrinsic uncertainty of the coefficient, not as a device for absorbing the laminar-versus-turbulent model-form uncertainty of the progression law: that model-form uncertainty is, in the author's parameterization, nominally carried by Sellmeijer's critical-head model factor $m_p \sim \mathrm{Lognormal}(1, 0.12)$ \parencite{pol_sie_2024}, and conflating the two would double-count it. The prototype-scale turbulent-transition evidence of \textcite{Okamura2022} motivates regarding the coefficient as genuinely uncertain at field scale, not as evidence that $C_e$ should stand in for the flow-regime model error. Second, because this engine integrates the pipe-growth rate over time, its natural anchor is the time-development target ($\approx 0.016$); adopting the higher mean-rate field value ($0.055$) is a deliberate, author-endorsed conservatism, and the difference is quantified rather than asserted. A propagation study over the two anchors (holding the remaining six parameters and the seepage length fixed) confirms that the lab-versus-field difference acts almost entirely through the \emph{mean}: the field prior raises the transient conditional failure probability by a factor of three to six at the fragility shoulder and two to three near the transition relative to the lab anchor, while the static branch, having no $C_e$ dependence, is unaffected; tightening only the coefficient of variation at fixed mean moves the result by at most ten per cent. The lab anchor $\mathrm{Lognormal}(0.016,\,0.48)$ is accordingly carried as the lower-bound robustness sensitivity in Chapter~\ref{chap: Discussion}, isolating the sensitivity of the calibrated result to the erosion-coefficient anchor.

\subsubsection*{The KP~63.4 outlier.}
KP~63.4 differs structurally from the other four sections in that its blanket is collapsed into a single gravelly foundation unit so that $k_{bl}$ is undefined, its aquifer conductivity is roughly one and a half orders of magnitude lower, it carries the thinnest blanket, and it is loaded by the shorter Shikaribetsu-referenced design event rather than the common Obihiro event. It is consequently treated as a distinct section. The decision as to whether it is pooled with the remaining four sections into a common fragility population is governed by the following pre-registered criterion: KP~63.4 is included in the pooled fragility population if its Phase~1 single-section conditional failure probability falls within the bootstrap confidence band of the four-section pooled curve; otherwise it is reported separately as a structurally anomalous section and its individual fragility is carried forward into Phase~2 and Phase~3 independently. This criterion is evaluated on the Phase~1 results; the outcome is not anticipated here.

\subsubsection*{Spatial extrapolation and reinforcement allocation.}
The priors of Tables~\ref{tab:priors_means} and~\ref{tab:priors_muln} are anchored exclusively to the five OYO cross-sections; no borehole or laboratory data from this dataset exist for the lower Tokachi reach or for the entire Satsunai left bank. These two sub-reaches, which together constitute a substantial portion of the study scope defined in Chapter~\ref{chap: Introduction}, must therefore carry interpolated prior means derived from the five measured sections, combined with deliberately inflated coefficients of variation that reflect the absence of direct geotechnical evidence. The inflated values are not yet assigned in the tables above and are determined as part of the Phase~1 spatial discretisation; they are carried forward as a primary source of epistemic uncertainty in the Phase~1 fragility generation and the Phase~2 calibration. The classification of each 200~metre node as drain-equipped, berm-only, or unreinforced, and the post-remediation seepage length $L$ for the reinforced segments, likewise depend on the current-state cross-sections that lie outside the OYO set \parencite{fukuda_2025_internal}. These dependencies are carried forward as explicit sources of epistemic uncertainty into the Phase~1 fragility generation and the Phase~2 calibration.

\section{The d4PDF Climate Ensemble: Structure, Downscaling, and Hydrograph Extraction}
\label{sec: The d4PDF Climate Ensemble: Structure, Downscaling, and Hydrograph Extraction}

The transient hydraulic loading that drives every limit state in this study is the river stage hydrograph $h(t)$ at each evaluation cross-section. Rather than adopting a single deterministic design hydrograph of the kind that underpins the historical safety assessments of the Tokachi basin \parencite{oyo_1999}, this study inherits a large-ensemble representation of the hydraulic forcing derived from the database for Policy Decision-Making for Future Climate Change (d4PDF). The use of this ensemble serves two distinct purposes. First, it replaces the limited instrumental record, typically several decades to a century, with thousands of years of physically consistent extreme-event realizations, allowing the spatiotemporal variability of flood loading to be treated as a quantifiable distribution rather than a single fixed scenario \parencite{mizuta_2017}. Second, and critically for the integration performed in Phase~3, the hydrograph ensemble used here is the same ensemble, processed through the same downscaling and runoff chain, that underlies the surface-failure fragility curves of \textcite{uemura_phd_2025}. This shared provenance ensures that the backward erosion piping fragility curves and the imported overflow and scour curves operate under identical hydraulic boundary conditions, which is a precondition for the series-system integration described in Chapter~\ref{chap: Methodology}.

\subsection{Structure of the d4PDF Ensemble and Dynamical Downscaling}
\label{subsec: Structure of the d4PDF Ensemble and Dynamical Downscaling}

The d4PDF database provides a very large ensemble of climate simulations produced by a 20~km regional atmospheric model, hereafter RCM20 \parencite{mizuta_2017}. It comprises a past experiment of 60 years (1951 to 2010) across 50 ensemble members, totalling 3{,}000 years, and a future experiment representing a 4~K global warming level (broadly equivalent to the RCP8.5 scenario toward the end of the twenty-first century) totalling 5{,}400 years. The future experiment is organized as 60 years across six prescribed sea-surface-temperature patterns, each with 15 members, yielding 90 member-equivalents of 60 years each. This distinction in internal structure, 50 flat members in the past experiment against six SST patterns of 15 members in the future experiment, is preserved in the present study's metadata so that the provenance of each hydrograph realization remains traceable.

The native 20~km resolution of RCM20 is too coarse to resolve the orographically controlled rainfall of the Tokachi basin and the flashy runoff response of its high-gradient tributaries. The ensemble was therefore dynamically downscaled to a 5~km regional model, hereafter RCM5, over a domain of 800~km by 800~km centred on 142.5\textdegree{}E, 42.75\textdegree{}N, with the domain extended southward over the Pacific to capture the development of approaching typhoons \parencite{yamada_2018}. Both the parent and the downscaled models employ the Meteorological Research Institute non-hydrostatic regional climate model. To make the downscaling of so large an ensemble computationally tractable, the simulation for each event was confined to a 15-day window bracketing the date of the annual-maximum basin-averaged rainfall, comprising a 5-day spin-up followed by the 5 days preceding and the 5 days following the event peak \parencite{yamada_2018}. The annual-maximum rainfall dates required to define these windows were identified from the RCM20 basin-averaged rainfall series, weighted by the fractional area of each climate-model grid cell intersecting the target catchment.

\subsection{From Precipitation to River Discharge}
\label{subsec: From Precipitation to River Discharge}

The downscaled RCM5 precipitation fields were converted to river discharge hydrographs using a storage-function runoff model for the Tokachi catchment \parencite{masuya2019spatiotemporal}. A central difficulty in applying a single calibrated runoff model across an ensemble spanning several orders of magnitude in rainfall is that parameters optimized to a large historical flood systematically overestimate runoff for the many smaller events in the ensemble. To address this, \textcite{uemura_phd_2025} adopted a rainfall-dependent parameterization in which the storage coefficient is expressed as a function of the 72-hour basin-averaged rainfall, calibrated against five reproduction events anchored on the September 2011 (H23.9) and August 2016 (H28.8) floods, the latter being the largest rainfall on record for the basin. This rainfall-dependent scheme reproduced the observed frequency distribution of annual-maximum peak discharge substantially better than a fixed-parameter calibration, recovering a median peak discharge close to the observed value at the Obihiro reference gauge.

Applied across the full ensemble, this runoff chain yields the distribution of flood discharge under both climate states. At the 1-in-150-year rainfall level, the median peak discharge at the Obihiro gauge increases by a factor of approximately 1.66 under the 4~K warming scenario relative to the past experiment, and a pronounced spread in peak discharge persists even at a fixed rainfall total, confirming that the temporal and spatial concentration of rainfall, and not merely its magnitude, governs the resulting flood loading \parencite{uemura_phd_2025}. This sensitivity to hydrograph shape rather than peak alone is precisely the feature that motivates the time-dependent treatment of backward erosion piping developed in this thesis.

\subsection{From Discharge to Transient River Stage}
\label{subsec: From Discharge to Transient River Stage}

The discharge hydrographs were converted to time series of river stage at each evaluation cross-section through a station-specific stage-discharge rating relation, with rating coefficients established from non-uniform flow computations performed at several discharge scales along the study reaches \parencite{uemura_phd_2025}. The channel geometry underlying these computations was taken from periodic cross-section survey data acquired in 2020 and 2021 for the Tokachi River and in 2022 for the Satsunai River, with channel roughness set by reproduction of recent major floods. The resulting stage hydrographs preserve the full temporal structure of each event, including the rising and receding limbs and, for compound events, the sequence of successive peaks. This continuity is essential to the present study: the time-stepping solver of Chapter~\ref{chap: Methodology} ingests the complete multi-peak series rather than a single peak value, and the cumulative-memory formulation depends on the inter-peak recession behaviour that only a continuous hydrograph can represent. The August 2016 consecutive-typhoon sequence, which produced the highest stage on record at the Obihiro gauge, reaching a level only 0.19~m below the design high-water level, is retained in its full transient form and serves as the empirical survival constraint for the Bayesian calibration of Phase~2.

\subsection{Hydrograph Extraction and Scenario Definition for the Present Study}
\label{subsec: Hydrograph Extraction and Scenario Definition for the Present Study}

For the reliability analysis, hydrographs are extracted at the 200-metre evaluation nodes spanning the Tokachi right bank from KP53.8 to KP62.8 and the Satsunai left bank from KP3.2 to KP16.6, the reaches fronting the Obihiro urban centre. Each extracted hydrograph carries the continuous stage series $h(t)$ together with its scenario tag, designating membership of either the 3{,}000-year past experiment or the 5{,}400-year 4~K warming experiment, and its native temporal resolution, which is retained so that the adequacy of the integration timestep against the steep rising limbs of the flashy peaks can be verified during the numerical analysis. The two ensembles are processed identically and differ only in their climate forcing, so that the difference between the resulting fragility curves isolates the climate-change signal, which is the object of the sensitivity assessment in Chapter~\ref{chap: Methodology}.

Two features of this dataset warrant emphasis for the interpretation of later results. First, because the runoff and rating chain was calibrated and constructed by \textcite{uemura_phd_2025} for the same evaluation nodes used here, the hydraulic loading is internally consistent across all three failure mechanisms integrated in Phase~3, and no additional hydraulic translation between datasets is required. Second, the design-scale discharges that anchor the conditioning range, namely 7{,}600~m\textsuperscript{3}/s for the Tokachi River and 2{,}900~m\textsuperscript{3}/s for the Satsunai River under the River Improvement Basic Policy, fall well within the envelope of the warming-scenario ensemble, confirming that the d4PDF ensemble spans the full range of loadings relevant to the reliability assessment without extrapolation beyond the simulated data.

\section{Pre-Calculated Surface Failure Fragility Curves from Uemura (2025)} \label{sec: Pre-Calculated Surface Failure Fragility Curves from Uemura (2025)}

To generate a comprehensive multi-mechanism risk profile without redundant computational overhead, pre-calculated conditional failure probability fragility curves for overflow and fluvial scour are imported directly from \textcite{uemura_phd_2025} and \textcite{uemura_wp2_2024}. These curves express the conditional probability of failure as a function of water level, $P_{f,i}(h)$, for each 200-metre evaluation segment and each d4PDF scenario ensemble. Because they were generated using the same d4PDF transient hydrographs and the same spatial discretization as this study, all three failure mechanism assessments share fully consistent hydraulic boundary conditions and spatial nodes, a prerequisite for the series-system integration of Phase~3.

The overflow fragility is based on cumulative damage theory applied to the time-integrated overtopping flow velocity \parencite{dean_2010}, while the fluvial scour fragility is based on a time-integrated bed shear stress exceedance model following USACE practice. Both account for relevant geotechnical and hydraulic uncertainties at the segment level. These curves are treated as pre-computed fixed inputs to the joint-probability integration; the fluvial scour model is applied in dimensionally consistent SI units, under which its contribution is negligible for these gravel levees, while the overflow curves are used as received.

\section{The Length Effect and Spatial Autocorrelation of Governing Parameters} \label{sec: The Length Effect and Spatial Autocorrelation of Governing Parameters}

Piping is a weakest-link failure mechanism: the probability of encountering a locally critical condition a thin blanket zone, an elevated hydraulic conductivity pocket, or a locally cohesionless region at the exit point increases monotonically with the physical length of the levee segment being assessed. Evaluating only a single representative cross-section therefore systematically underestimates the failure probability of a finite-length segment \parencite{kanning_2012}. This spatial scale effect is accounted for by scaling per-cross-section conditional failure probabilities to the 200-metre segment level using the spatial autocorrelation length of the governing parameters $D_{bl}$ and $k_\mathrm{aquifer}$, following the approach of \textcite{hoffmans_2014}.

For the portions of the study reach covered by the OYO (1999) borehole survey specifically KP56.0 to KP66.2 on the Tokachi right bank and KP3.4 to KP7.0 on the Satsunai left bank the spatial autocorrelation length is estimated empirically from the variability of these parameters across the surveyed cross-sections, providing a site-specific basis for the correction. The continuous soil layer longitudinal profile (土層縦断図) in the report appendix \parencite{oyo_1999} is particularly valuable for this characterisation, as it documents the along-levee variation of the $A_c$ layer thickness at a resolution substantially finer than the individual cross-section spacing.

For the data-gap segments (KP53.8--56.0 on the Tokachi right bank and KP7.0--16.6 on the Satsunai left bank), where continuous borehole coverage is absent, autocorrelation lengths are assigned from representative values in the literature for comparable alluvial levee foundation systems \parencite{kanning_2012, hoffmans_2014}, and the resulting Length Effect corrections carry correspondingly higher uncertainty. This is explicitly acknowledged in the uncertainty analysis of Chapter~\ref{chap: Results: Subsurface Piping Assessment}.
